\chapter{Tensorization}\label{parTensorization}%
\addcontentsline{itc}{chapter}{\protect\numberline{\thechapter}Tensorization}

\section{Subjective correlation}\label{parSubjectiveCorrelation}

In this chapter we will need more advanced definitions for decorrelation.

\begin{Def}
Let~$X$, $Y$ and~$Z$ be random variables. For~$\epsilon\geq 0$,
one says that $X$ and~$Y$ are \emph{subjectively $\epsilon$-decorrelated w.r.t.~$Z$}
(or \emph{$\epsilon$-decorrelated seen from $Z$})
if $X$ and~$Y$ are $\epsilon$-decorrelated under the law $\Law(X,Y|{Z=z})$
for $\Law(Z)$-almost-all~$z$\footnote{The conditional laws $\Law(\Bcdot|{Z=z})$
are only defined up to~$\Law(Z)$-a.e.\ equality, whence the need to specify
``for $\Law(Z)$-almost-all~$z$''.}.

The smallest $\epsilon$ such that $X$ and~$Y$ are $\epsilon$-decorrelated
seen from $Z$ will be called the \emph{subjective correlation level
between~$X$ and~$Y$ w.r.t.~$Z$} (or \emph{correlation level between~$X$ and~$Y$ seen from $Z$});
we denote it ${\{X:Y\}}_{Z}$.
\end{Def}

In \S~\ref{parDefinition}, we had given the definitions
in terms of $\sigma$-algebras rather than random variables.
Of course there is also a $\sigma$-algebra definition for subjective correlation,
though I find it harder to understand:
%
\begin{Def}\label{def5436}
Let~$\mcal{F}$, $\mcal{G}$ and~$\mcal{H}$ be $\sigma$-algebras.
For~$\epsilon \in [0,1]$, the expression
``$\{\mcal{F}:\mcal{G}\}_{\mcal{H}} \leq \epsilon$'' means that
for all~$f \in \ldb(\mcal{F}\vee\mcal{H})$ and all~$g \in \ldb(\mcal{G}\vee\mcal{H})$
satisfying $\EE[f|\mcal{H}] \equiv 0$, resp.\ $\EE[g|\mcal{H}] \equiv 0$,
one has:
\[ |\EE[fg]| \leq \epsilon \ecty(f) \ecty(g) .\]
\end{Def}

We let the reader check that with that definition,
for~$X$, $Y$ and~$Z$ random variables,
${\{X:Y\}}_{Z} = \{\sigma(X):\sigma(Y)\}_{\sigma(Z)}$.

\begin{Rmk}
The ordinary correlation can be seen as a particular case of subjective correlation,
since $\{\mcal{F}:\mcal{G}\} = \{\mcal{F}:\mcal{G}\}_{\mcal{O}}$ for $\mcal{O} = \{\emptyset,\Omega\}$
the trivial $\sigma$-field.
\end{Rmk}

\begin{Rmk}\label{r4670}
Warning! Writing that ${\{X:Y\}}_Z \leq \epsilon$ does \emph{not} imply
that for all subset $C$ of the range of~$Z$,
$X$ and~$Y$ are $\epsilon$-decorrelated under~$\Law(X,Y|\ab Z\in C)$:
see Examples~\ref{xpl3309} and~\ref{xpl3450} below.
\end{Rmk}

\begin{Rmk}
Warning again! There is no general inequality between~%
${\{X:Y\}}$ and~${\{X:Y\}}_Z$: see Examples~\ref{xpl3308} and~\ref{xpl3309} below.
\end{Rmk}

\begin{Xpl}
Let~$f \colon \RR \to \RR_+$ be a nonnegative continuous function
with $\int_{\RR} f(x)\,\dx{x} = 1$ and let~$(X,Y,Z)$ be a variable
on~$\RR^3$ with density
\[ \dx\Pr\big[(X,Y,Z)=(x,y,z)\big] = \frac{1}{2\pi} f(z)
\exp\Big( \sinh z\cdot xy - {\textstyle\frac{1}{2}} \cosh z\cdot(x^2+y^2) \Big) \,\dx{x}\dx{y}\dx{z} .\]
Then, conditionally to ``$Z=z$'', $(X,Y)$ is a Gaussian vector
with $\Var(X) = \Var(Y) = \cosh z$ and $\Cov(X,Y) = \sinh z$,
so by Theorem~\ref{pro1857}, under the law $\Pr[\Bcdot|Z=z]$ one has ${\{X:Y\}} = |\tanh z|$.
Consequently ${\{X:Y\}}_Z = \sup\{ |\tanh z| \colon f(z)>0 \}$.
\end{Xpl}

The three following examples show that subjective correlation may behave rather wildly,
especially when one changes the $\sigma$-field of reference:
%
\begin{Xpl}\label{xpl3308}
Let~$X$ and~$Y$ be independent variables with uniform law on~$\RR/\ZZ$
and let~$Z=X+Y$; then ${\{X:Y\}}_{Z} = 1$: under~$\Pr[\Bcdot|\ab Z=z]$ indeed
$Y$ is $X$-measurable (and not constant), since $Y\equiv z-X$.
\end{Xpl}

\begin{Xpl}\label{xpl3309}
Let~$\alpha,\beta,\gamma$ be three independent random variables
uniform on~$\{0,1\}$; define ${X=(\gamma,\alpha)}, \ab Y=(\gamma,\beta)$ and $Z=\gamma$.
Then, conditionally to ``$Z=0$'', $X$ and~$Y$ are independent
with common law uniform on~$\{(0,0),(0,1)\}$,
and similarly $X$ and~$Y$ are independent conditionally to ``$Z=1$'',
so ${\{X:Y\}}_Z=0$.
Yet $X$ and~$Y$ are not independent
since the events ``$X\in\{(0,0),(0,1)\}$'' and ``$Y\in\{(0,0),(0,1)\}$'',
which are non-trivial under~$\Pr$, are equivalent,
so that ${\{X:Y\}}=1$.
\end{Xpl}

\begin{Xpl}\label{xpl3450}
Let $X=(X_1,X_2)$ and $Y=(Y_1,Y_2)$ be independent with uniform laws
on~$\{0,1\}^2$ and define $Z=(X_1,Y_1)$;
then one easily checks that ${\{X:Y\}}_Z=0$.
Now let~$Z'=\1{X_1=Y_1}$, which is $Z$-measurable;
one has ${\{X:Y\}}_{Z'}=1$ since, for instance,
under ``$Z'=1$'' the events ``$X_1=0$'' and ``$Y_1=0$''
are non-trivial and equivalent.
\end{Xpl}

Now we define a more restrictive concept of subjective correlation.
%
\begin{Def}
A \emph{$\sigma$-metalgebra} $\mcal{M}$ is a set $\{\mcal{H}\colon \mcal{H}\in\mcal{M}\}$
of $\sigma$-algebras which is stable under the ``$\bigvee$'' operator,
i.e.\ such that for any~$\mcal{M}'\subset\mcal{M}$, $\bigvee_{\mcal{H}\in\mcal{M}'}\mcal{H}\in\mcal{M}$.
\end{Def}

One can speak of the `$\sigma$-metalgebra spanned by some set of $\sigma$-algebras',
as states the following immediate proposition:
%
\begin{Pro}\label{pro0539}
If $(\mcal{H}_k)_{k\in K}$ is a set of $\sigma$-algebras, then there is a smallest
$\sigma$-metalgebra containing all the~$\mcal{H}_k$, which is
\[ \mcal{M} = \Big\{ \bigvee_{k\in K'} \mcal{H}_k \ ;\ K'\subset K \Big\} .\]
\end{Pro}

When one deals with random variables rather than $\sigma$-algebras,
one has the following variant of Proposition~\ref{pro0539}:
%
\begin{Pro}\label{cor0780}
Let~$(Z_k)_{k\in K}$ be a set of random variables,
then the $\sigma$-metalgebra spanned by~$\{\sigma(Z_k) \colon k\in K\}$
is $\{\sigma(\vec{Z}_{K'})\colon K'\subset K\}$.
\end{Pro}

\begin{Def}
Let~$\mcal{F}$ and~$\mcal{G}$ be $\sigma$-algebras and $\mcal{M}$ be a $\sigma$-metalgebra.
We define the \emph{correlation between~$\mcal{F}$ and~$\mcal{G}$ seen from $\mcal{M}$} by:
\[ \{\mcal{F}:\mcal{G}\}_{\mcal{M}} = \sup_{\mcal{H}\in\mcal{M}} \{\mcal{F}:\mcal{G}\}_{\mcal{H}} .\]
\end{Def}

\begin{Rmk}
Speaking in terms of random variables, if $X$, $Y$ and~$(Z_k)_{k\in K}$ are variables,
denoting by~$\mcal{M}$ the $\sigma$-metalgebra spanned by the~$Z_k$,
then ${\{X:Y\}}_{\mcal{M}}$ is the supremum%
\footnote{More precisely it is a \emph{true} supremum (over $K'$)
of \emph{essential} suprema (over $\vec{z}_{K'}$).}
of the~${\{X:Y\}}$ when taken under all the laws
of kind $\Pr[\Bcdot|\vec{Z}_{K'}=\vec{z}_{K'}]$ for~$K'$ a subset of~$K$
and~$z_k, k\in K'$ elements of the respective ranges of the~$Z_k$.
\end{Rmk}

Finally, the following proposition gathers some easy properties
of relative correlation w.r.t.\ a $\sigma$-metalgebra:
%
\begin{Pro}\strut
\begin{ienumerate}
\item Call~$\mcal{M}_{\emptyset}$ the trivial $\sigma$-metalgebra, that is,
$\mcal{M}_{\emptyset} = \{\mcal{O}\}$; then for all $\sigma$-algebras~$\mcal{F}$ and~$\mcal{G}$,
$\{\mcal{F}:\mcal{G}\} = \{\mcal{F}:\mcal{G}\}_{\mcal{M}_{\emptyset}}$.
\item If $\mcal{M} \subset \mcal{M}'$, then $\{\mcal{F}:\mcal{G}\}_{\mcal{M}} \leq \{\mcal{F}:\mcal{G}\}_{\mcal{M}'}$.
\item Let~$\mcal{F}$ and~$\mcal{G}$ be $\sigma$-algebras,
let~$\mcal{M}$ be a $\sigma$-metalgebra, and call~$\tilde{\mcal{M}}$ the $\sigma$-metalgebra
spanned by~$\mcal{M}$, $\mcal{F}$ and~$\mcal{G}$; then
$\{\mcal{F}:\mcal{G}\}_{\mcal{M}} = \{\mcal{F}:\mcal{G}\}_{\tilde{\mcal{M}}}$.
\end{ienumerate}
\end{Pro}

\begin{Def}\label{def3128}
In the sequel, the probabilistic systems which we shall consider
will often be made of some `elementary' variables, say $(X_i)_{i\in I}$.
In this case, the so-called \emph{natural $\sigma$-metalgebra} of the system
will mean the $\sigma$-metalgebra spanned by the~$X_i$.
\end{Def}




\section{Simple tensorization}\label{parSimpleTensorization}

Now we turn to tensorization.
First let us deal with `simple' tensorization,
by which I mean that tensorization is performed on only one variable.
The main result of this section will be the `$N$~against~$1$' theorem
(Theorem~\ref{thm5750}).

The problem considered is the following:
Let~$I$ be a set and~$(X_i)_{i\in I}, Y$ be random variables;
call~$\mcal{M}$ the natural $\sigma$-metalgebra of this system, that is,
the $\sigma$-metalgebra spanned by the~$X_i$ and~$Y$ (cf.\ Definition~\ref{def3128}).
Suppose we have bounds ${\{X_i:Y\}_{\mcal{M}}} \leq \epsilon_i$ for all~$i$;
the question is, can we deduce from them a bound on~$\{\vec{X}_I:Y\}$?
We shall prove that the answer is ``yes'',
and moreover the bound~(\ref{for5719}) we will give is optimal in some way
(see \S~\ref{parOptimality}).

For pedagogical purpose,
let us first state and prove a weaker but easier proposition:
%
\begin{Pro}\label{pro5527}
With the notation above,
\[\label{for5528} \big\{\vec{X}_I:Y\big\} \leq \sqrt{\sum_{i\in I} \epsilon_i^2} .\]
\end{Pro}

\begin{proof}
By Proposition~\ref{pro6559} we may assume $I=\{1,\ldots,N\}$.
Let~$f$ and~$g$ be centered $L^2$ $\vec{X}$-measurable,
resp.\ $Y$-measurable, functions; our goal is to bound $|\EE[fg]|$.

For all~$i\in\{0,\ldots,N\}$, denote
\[ \mcal{F}_i = \sigma(X_1,\ldots,X_i) ,\]
and for all~$i\in\{1,\ldots,N\}$,
\[ f_i = f^{\mcal{F}_i} - \EE[f|\mcal{F}_{i-1}] .\]
%
Then $f = \sum_i f_i$, where each~$f_i$ is $\mcal{F}_i$-measurable
and centered w.r.t.~$\mcal{F}_{i-1}$ (i.e., $\EE[f_i|\mcal{F}_{i-1}] \equiv 0$).
Consequently, for all~$i_0<i_1$ one has $\EE[f_{i_0}f_{i_1}] = 0$
(since $f_{i_0}$ is $\mcal{F}_{i_0}$-measurable while
$f_{i_1}$ is centered w.r.t.~$\mcal{F}_{i_1-1} \supset \mcal{F}_{i_0}$)
and thus when one expands $\Var(f) = \EE\big[\big(\sum_i f_i\big)^2\big]$
all the non-diagonal terms vanish, yielding:
\[\label{f9207} \Var(f) = \sum_{i=1}^{N} \Var(f_i) .\]

Now, the decomposition ``$f = \sum_i f_i$'' yields
\[ \EE[fg] = \sum_{i=1}^{N} \EE[f_ig] ,\]
so let us bound the~$|\EE[f_ig]|$.
The law of total expectation gives:
\[ \EE[f_ig] =
\int \EE\big[f_ig\big|X_1=x_1,\ldots,X_{i-1}=x_{i-1}\big] \,\dx\Pr[x_1,\ldots,x_{i-1}] .\]
But under~$\dx\Pr[\Bcdot|\ab x_1,\ldots,x_{i-1}]$,
$f_i$ is $X_i$-measurable and centered while $g$ is $Y$-mea\-su\-ra\-ble,
moreover under this law
$\{X_i:Y\} \leq \{X_i:Y\}_{\mcal{F}_{i-1}} \leq \{X_i:Y\}_{\mcal{M}} \leq \epsilon_i$, so:
\[ \big|\EE\big[f_ig\big|Y_1=y_1,\ldots,Y_{i-1}=y_{i-1}\big]\big| \leq \epsilon_i
\ecty\big(f_i\big|x_1,\ldots,x_{i-1}\big) \ecty\big(g\big|x_1,\ldots,x_{i-1}\big) .\]
%
Using the bound $\ecty(h) \leq \sqrt{\EE[h^2]}$, it follows that:
\begin{multline}\label{cal4816}
\big|\EE[f_ig]\big| \leq \epsilon_i
\int \sqrt{\EE[f_i^2|x_1,\ldots,x_{i-1}]} \sqrt{\EE[g^2|x_1,\ldots,x_{i-1}]}
\,\dx\Pr[x_1,\ldots,x_{i-1}] \\
\footrel{\text{CS}}{\leq} \epsilon_i \sqrt{\int\EE[f_i^2|x_1,\ldots,x_{i-1}]\,\dx\Pr[x_1,\ldots,x_{i-1}]}\sqrt{\textit{the same for~$g$}}
= \epsilon_i \ecty(f) \ecty(g_i) .
\end{multline}
So, summing~(\ref{cal4816}) for all~$i$:
\[\label{f3109} |\EE[fg]| \leq \sum_{i=1}^N \epsilon_i \ecty(f_i) \ecty(g)
\footrel{\text{CS}}{\leq} \sqrt{\sum_i \epsilon_i^2} \sqrt{\sum_i \Var(f_i)} \, \ecty(g)
= \sqrt{\sum_i \epsilon_i^2} \, \ecty(f) \ecty(g) .\]
Since (\ref{f3109}) is true for all~$f,g$, (\ref{for5528}) is proved.
\end{proof}

\label{lbl91}It is striking in Proposition~\ref{pro5527}
that the right-hand side of~(\ref{for5528})
may be greater than $1$, which is never the case for a correlation level.
Actually there is some `loss of optimality' in the proof
of the proposition when we bound above $\Var(f_i|\mcal{F}_{i-1})$
by~$\EE[f_i^2|\mcal{F}_{i-1}]$, since $\EE[f_i^2|\mcal{F}_{i-1}] - \Var(f_i|\mcal{F}_{i-1}) = \EE[f_i|\mcal{F}_{i-1}]^2$
may be different to~$0$. We will use a technique for `recycling' that loss
to get the following result, which \S~\ref{parOptimality} shall prove to be optimal:
%
\begin{Thm}[`$N$~against~$1$' theorem]\label{thm5750}
Take the same hypotheses as in Proposition~\ref{pro5527}:
$\forall i\in I \quad \{X_i:Y\}_{\mcal{M}} \leq \epsilon_i$, where $\mcal{M}$
is the natural $\sigma$-metalgebra of the system. Then:
\[\label{for5719} \big\{\vec{X}_I:Y\big\} \leq \sqrt{1-\prod_{i\in I}(1-\epsilon_i^2)} .\]
\end{Thm}

\begin{Rmk}
The right-hand side of~(\ref{for5719})
is the~$\bar\epsilon\in[0,1]$ characterized by
$1-\bar\epsilon^2 = \prod_i (1-\epsilon_i)^2$.
\end{Rmk}

\begin{Rmk}
The right-hand side of~(\ref{for5719}) is bounded above by~$\sqrt{\sum_i \epsilon_i^2}$,
so Theorem~\ref{thm5750} gives back Proposition~\ref{pro5527} as a corollary.
\end{Rmk}
%
\begin{proof}
As in the proof of Proposition~\ref{pro5527},
let~$f$ and~$g$ be centered $L^2$ $\vec{X}$-measurable,
resp.\ $Y$-measurable, functions.
Assume $I=\{1,\ldots,N\}$;
denote $\mcal{F}_i \coloneqq \sigma(X_1,\ldots,X_i)$
and $f_i \coloneqq f^{\mcal{F}_i} - \EE[f|\mcal{F}_{i-1}]$.
Also denote, for~$i\in\{0,\ldots,N\}$,
\[ g^i \coloneqq g - \EE[g|\mcal{F}_i] .\]

As before, one has $\Var(f) = \sum_i\Var(f_i)$ and $\EE[fg] = \sum_{i=1}^N \EE[f_ig]$.
But $f_i$ is centered w.r.t.~$\mcal{F}_{i-1}$
while $(g - g^{i-1})$ is $\mcal{F}_{i-1}$-measurable,
so $\EE[f_ig] = \EE[f_ig^{i-1}]$.
Since, conditionally to~$\mcal{F}_{i-1}$, $f_i$ and~$g^{i-1}$ are both centered and resp.\
$X_i$- and $Y$-measurable, the fact that $\{X_i:Y\}_{\mcal{M}} \leq \epsilon_i$ implies,
by the same argument as in the previous proof,
that
\[ |\EE[f_ig^{i-1}]| \leq \epsilon_i \ecty(f_i) \ecty(g^{i-1}) .\]

Now, for~$i\in\{1,\ldots,N\}$, denote
\[ \bar{g}^i = \EE[g^{i-1}|\mcal{F}_i] .\]
Since $g^{i-1} = \bar{g}^i + g^i$,
where $\bar{g}^i$ is $\mcal{F}_i$-measurable while $g^i$ is centered w.r.t.~$\mcal{F}_i$, one has:
\[ \Var(g^i) = \Var(g^{i-1}) - \Var(\bar{g}^i) .\]
%
Then, the point consists in making the following observation:
for~$\Var(g^i)$ to be large (that is, close to~$\Var(g^{i-1})$),
$\Var(\bar{g}^i)$ has to be small. But in that case
$|\EE[f_ig]|$ shall be small:
one has indeed, since $f_i$ is $\mcal{F}_i$-measurable,
\[ |\EE[f_ig]| = |\EE[f_ig^{i-1}]|
= \big|\EE\big[f_i(g^{i-1})^{\mcal{F}_i}\big]\big| = \big|\EE\big[f_i\bar{g}^i\big]\big|
\footrel{\text{CS}}{\leq} \ecty(f_i) \ecty(\bar{g}^i) .\]

Let us sum up the relations obtained.
One has, for all~$i\in\{1,\ldots,N\}$:
\begin{eqnarray}
\label{f9041a} |\EE[f_ig]| & \leq & \epsilon_i \ecty(f_i) \ecty(g^{i-1}) ; \\
\label{f9041b} \ecty(g^i) & = & \sqrt{\ecty(g^{i-1})^2 - \ecty(\bar{g}^i)^2} ; \\
\label{f9041c} |\EE[f_ig]| & \leq & \ecty(f_i) \ecty(\bar{g}^i) .
\end{eqnarray}
%
Now define $\hat\epsilon_i \coloneqq |\EE[f_ig]| \div \ecty(f_i)\ecty(g^{i-1})$,
or $\hat\epsilon_i = 0$ if the right-hand side is $0 \div 0$.
Then (\ref{f9041a}) ensures that $\hat\epsilon_i \leq \epsilon_i$,
and (\ref{f9041c}) means that $\ecty(\bar{g}^i) \geq \hat\epsilon_i\ecty(g^{i-1})$,
so that (\ref{f9041b}) yields $\ecty(g^i) \leq \sqrt{1-\hat\epsilon_i^2} \*
\hspace{.056em} \ecty(g^{i-1})$.
Since $g^0 = g$, one has therefore by induction
$\ecty(g^i) \leq \prod_{i'=1}^{i-1} \ab \sqrt{1-\hat\epsilon_{i'}^2} \,\ecty(g)$,
so that the decomposition ``$\EE[fg]=\sum_i\EE[f_ig]$'' gives:
\[\label{f8485} |\EE[fg]| \leq \sum_{i=1}^N
\Big( \hat\epsilon_i \prod_{i'=1}^{i-1} \sqrt{1-\hat\epsilon_{i'}^2} \Big)
\ecty(f_i) \ecty(g) .\]
By the Cauchy--Schwarz inequality, it follows that:
\[\label{f9280} |\EE[fg]| \leq \sqrt{\sum_{i=1}^N \hat\epsilon_i^2
\prod_{i'=1}^{i-1} \big(1-\hat\epsilon_{i'}^2 \big)} \, \ecty(f)\ecty(g)
= \sqrt{1-\prod_{i=1}^N \big(1-\hat\epsilon_i^2\big)} \, \ecty(f)\ecty(g) .\]
%
Obviously the maximal value for the right-hand side of~(\ref{f9280})
is when $\hat\epsilon_i = \epsilon_i$ for all~$i$,
then yiel\mbox{di}ng~(\ref{for5719}).
\end{proof}

There is an alternative proof, which is less intuitive
but whose reasoning shall be used again in the proof of Theorem~\ref{thm6413}:
%
\begingroup\def\proofname{Alternative proof of Theorem~\ref{thm5750}}
\begin{proof}
We use the same notation as in the previous proof.
As $f$ is $\mcal{F}$-measurable, $\EE[fg] = \EE[fg^{\mcal{F}}]$,
so by the Cauchy--Schwarz inequality:
\[\label{f2641} |\EE[fg]| \leq \ecty(f) \ecty(g^{\mcal{F}}) .\]
Now, by associativity of variance $\ecty(g^{\mcal{F}}) = \sqrt{\Var(g)-\Var(g-g^{\mcal{F}})}$,
so by~(\ref{f2641}) it suffices to prove that
\[\label{for5211}
\Var\big(g-g^{\mcal{F}}) \geq \prod_{i=1}^N \big(1-\epsilon_i^2\big) \Var(g) .\]

With our notation, $g-g^{\mcal{F}} = g^N$ and $g = g^0$;
we will prove that for all~$i\in\{1,\ldots,N\}$,
\[\label{for0829} \Var(g^i) \geq (1-\epsilon_i^2) \Var(g^{i-1}) .\]
%
Since $g^{i-1}$ and~$g^i$ are centered w.r.t.~$\mcal{F}_{i-1}$,
one has
\[ \Var(g^{i-1}) = \int \Var(g^{i-1}|x_1,\ldots,x_{i-1}) \,\dx\Pr[x_1,\ldots,x_{i-1}] ,\]
with a similar decomposition for~$\Var(g^i)$,
so that it suffices to prove~(\ref{for0829}) conditionally to~$\mcal{F}_{i-1}$.

Conditionally to~$\mcal{F}_{i-1}$, $g^{i-1}$ is centered and $Y$-measurable.
Moreover, $g^i = g^{i-1} - (g^{i-1})^{\sigma(X_i)}$,
so by associativity of variance
$\Var(g^i) = \Var(g^{i-1}) - \Var\big( (g^{i-1})^{\sigma(X_i)} \big)$,
and therefore (\ref{for0829}) is equivalent to
\[ \Var\big( (g^{i-1})^{\sigma(X_i)} \big) \geq \epsilon_i^2 \Var(g^{i-1}) ,\]
which follows directly from the assumption ``$\{X_i:Y\}_{\mcal{F}_{i-1}} \leq \epsilon_i$''.
\end{proof}\endgroup




\section{Double tensorization}\label{parDoubleTensorization}

Simple tensorization as itself is already interesting
since it gives an $L^2$-type bound for the correlation between~$X$ and~$\vec{Y}$,
which is better than the $L^1$-type bounds
typically obtained by total variation methods.
Yet it does not exhaust
the full potential of Hilbertian correlations
concerning tensorization, since obviously
it does not contain results like independent tensorization
(cf.\ \S~\ref{parIndependentTensorization}).

The aim of this section is to get sharp tensorization results
where we perform tensorizing on \emph{both} sides,
without having to assume complete independence like in Theorem~\ref{pro1252}.
The price to pay is that the techniques involved,
though similar in their spirit, will be much more tricky,
moreover the bounds obtained will not be completely optimal
(see \S~\ref{parOptimality}).


\subsection{`\texorpdfstring{$N$}{N}~against~\texorpdfstring{$M$}{M}' tensorization}

The following theorem may be considered as the main result of this monograph.
As will be explained in \S~\ref{parAsymptoticOptimality},
it `contains' qualitatively all the other tensorization theorems
(i.e.\ Theorems~\ref{pro1252}, \ref{thm5750} and~\ref{thm0119}).

\begin{Thm}[`$N$~against~$M$' theorem]\label{thm6413}
Let~$I$ and~$J$ be sets, and let~$(X_i)_{i\in I}$ and~$(Y_j)_{j\in J}$ be random variables,
the $\sigma$-metalgebra they generate being denoted by~$\mcal{M}$.
Suppose for any~$i,j$, ${\{X_i:Y_j\}_{\mcal{M}}} \leq \epsilon_{ij}$ for some $\epsilon_{ij} \geq 0$,
and define the operator
\[ \begin{array}{rrcl} \boldepsilon \colon & L^2(J) & \to & L^2(I) \\
& (a_j)_{j\in J} & \mapsto & \big( \sum_{j\in J} \epsilon_{ij}a_j \big)_{i\in I} ,\end{array} \]
then:
\[\label{for3009} \big\{ \vec{X}_I : \vec{Y}_J \big\} \leq \VERT \boldepsilon \VERT \wedge 1 .\]
\end{Thm}

\begin{Rmk}
On~$(\RR_+)^{I\times J}$, $\VERT \boldepsilon \VERT$
is a nondecreasing function of each~$\epsilon_{ij}$.
\end{Rmk}

\begin{NOTA}
As the proof of Theorem~\ref{thm6413} is rather technical,
I found it useful to write down how it goes on a concrete example.
This is performed in Appendix~\ref{parIllustrationOfTheProofs},
which I suggest the reader to look at in parallel with the proof as a complement.
\end{NOTA}

To prove Theorem~\ref{thm6413},
we will need the following
%
\begin{Lem}\label{lem7265}
Let~$X_1,X_2,\ldots,X_N$ and~$Y$ be random variables,
call~$\mcal{M}$ their natural $\sigma$-metalgebra,
and assume that for all~$i\in\{1,\ldots,N\}$,
\[ \{X_i:Y\}_{\mcal{M}} \leq \epsilon_i .\]
Let~$f$ be an $\ldb(\vec{X})$ function.
For all~$0 \leq i \leq N$, denote $\mcal{F}_i \coloneqq \sigma(X_1,\ldots,X_i)$,
resp.\ $\mcal{F}^*_i \coloneqq \sigma(X_1,\ldots,\ab X_i,Y)$,
and for all~$1 \leq i \leq N$, define
\[ f_i \coloneqq f^{\mcal{F}_i} - \EE[f|\mcal{F}_{i-1}] ,\]
\[ \llap{\text{resp.}\quad} f^*_i \coloneqq f^{\mcal{F}^*_i} - \EE[f|\mcal{F}^*_{i-1}] ,\]
and denote by~$V_i$ and~$V^*_i$ their respective variances. Then,
for all~$1 \leq i \leq N$,
\[\label{for1302}
V^*_i \geq (1-\epsilon_i^2) V_i - 2\epsilon_i \sqrt{V_i} \Big( \sum_{i'>i} \epsilon_{i'} \sqrt{V_{i'}} \Big) .\]
\end{Lem}

\begin{proof}
For~$0\leq i\leq N$, define
\[ \tilde{f}_i \coloneqq f - \EE[f|\mcal{F}_i] ,\]
\[ \llap{\text{resp.}\quad} \tilde{f}^*_i \coloneqq f - \EE[f|\mcal{F}^*_i] ,\]
and call~$\tilde{V}_i$ and~$\tilde{V}^*_i$ their respective variances.
One has $\tilde{f}_i = \sum_{i'>i} f_{i'}$, resp.\ $\tilde{f}^*_i = \sum_{i'>i} f^*_{i'}$.
Moreover, by the same argument as in the proof of Proposition~\ref{pro5527},
all the~$f_i$ are orthogonal (that is, ${i_0\neq i_1} {\enskip\Rightarrow}\enskip \EE[f_{i_0}f_{i_1}] = 0$),
thus
\[ \tilde{V}_i = \sum_{i'>i} V_{i'} ;\]
similarly,
\[ \tilde{V}^*_i = \sum_{i'>i} V^*_{i'} .\]

In a first step, we observe that for all~$i$,
$\tilde{f}_i - \tilde{f}^*_i = (f-f^{\mcal{F}_i}) - (f-f^{\mcal{F}^*_i})
= f^{\mcal{F}^*_i} - f^{\mcal{F}_i} = {f^{\mcal{F}^*_i} - (f^{\mcal{F}_i})^{\mcal{F}^*_i}} \ab = (f - f^{\mcal{F}_i})^{\mcal{F}^*_i} = (\tilde{f}_i)^{\mcal{F}^*_i}$,
which by associativity of variance yields the following
%
\begin{Clm}\label{clm1908}
\[\label{for4854} \tilde{V}_i - \tilde{V}^*_i = \Var\big((\tilde{f}_i)^{\mcal{F}^*_i}\big).\]
\end{Clm}

Now, the following claim will be the main tool for proving the lemma:
%
\begin{Clm}\label{clm1959}
For all~$1 \leq i \leq N$,
\[\label{for1960}
\tilde{V}_{i-1} - \tilde{V}^*_{i-1} \leq \Big( \epsilon_i \sqrt{V_i} + \sqrt{\tilde{V}_i-\tilde{V}^*_i} \Big)^{\!2\,}
\footnotemark .\]
\end{Clm}
%
\footnotetext{Taking the square root of~$(\tilde{V}_i-\tilde{V}^*_i)$
is allowed, since that quantity is nonnegative by Claim~\ref{clm1908}.}

Admit temporarily Claim~\ref{clm1959}. Since $\tilde{V}_N= \tilde{V}^*_N = 0$,
(\ref{for1960}) applied with $i=N$ gives $\tilde{V}_{N-1} - \tilde{V}^*_{N-1} \leq \epsilon_N^2 V_N$,
which in turn we can use in~(\ref{for1960}) with $i=N-1$, and so on,
to finally prove by finite (decreasing) induction that, for all~$i$,
\[\label{for1291}
\tilde{V}_i - \tilde{V}^*_i \leq \Big( \sum_{i'>i}\epsilon_{i'}\sqrt{V_{i'}} \Big)^2 .\]
Now to get~(\ref{for1302}),
we note that $V_i = \tilde{V}_{i-1} - \tilde{V}_i$,
resp.\ $V^*_i = \tilde{V}^*_{i-1} - \tilde{V}^*_i$,
so, using successively the inequalities~(\ref{for1960}) and~(\ref{for1291}),
\begin{multline}
V_i - V^*_i = \big( \tilde{V}_{i-1} - \tilde{V}^*_{i-1} \big)
- \big( \tilde{V}_i - \tilde{V}^*_i \big) \\
\leq \Big( \epsilon_i \sqrt{V_i} + \sqrt{\tilde{V}_i - \tilde{V}^*_i} \Big)^2
- (\tilde{V}_i - \tilde{V}^*_i)
= \epsilon_i^2 V_i + 2 \epsilon_i\sqrt{V_i} \sqrt{\tilde{V}_i - \tilde{V}^*_i} \\
\leq \epsilon_i^2 V_i + 2 \epsilon_i\sqrt{V_i} \Big( \sum_{i'>i}\epsilon_{i'}\sqrt{V_{i'}} \Big),
\end{multline}
which is equivalent to (\ref{for1302}).
\end{proof}

\begingroup\def\proofname{Proof of Claim~\ref{clm1959}}\begin{proof}
Thanks to Claim~\ref{clm1908}, what we have to prove is:
\[\label{for3998}
\Var\big((\tilde{f}_{i-1})^{\mcal{F}^*_{i-1}}\big) \leq \Big( \epsilon_i \sqrt{V_i} + \sqrt{\tilde{V}_i-\tilde{V}^*_i} \Big)^{\!2\,} .\]
By the definition of conditional expectation
and the equality case in the Cauchy--Schwarz inequality,
(\ref{for3998}) is equivalent to saying that for all~$\ldb(\mcal{F}^*_{i-1})$ function~$g$,
\[\label{for3331} \big| \EE \big[ \tilde{f}_{i-1} g \big] \big| \leq
\Big( \epsilon_i \sqrt{V_i} + \sqrt{\tilde{V}_i-\tilde{V}^*_i} \Big) \ecty(g) .\]
So let~$g$ be a centered $L^2$ $\mcal{F}^*_{i-1}$-measurable real function.
Since $\tilde{f}_{i-1} = f_i + \tilde{f}_i$,
$\EE[\tilde{f}_{i-1}g] = \EE[f_ig] + \EE[\tilde{f}_ig]$,
which two terms we shall bound separately.

For the first term, under~$\Pr[\Bcdot|\ab \mcal{F}_{i-1}]$,
$f_i$ is centered and only depends on~$X_i$, and~$g$ only depends on~$Y$.
Since $\{ X_i : Y \}_{\mcal{F}_{i-1}} \leq {\{ X_i : Y \}_{\mcal{M}}} \leq \epsilon_i$,
it follows that
\[ \big| \EE[f_ig|\mcal{F}_{i-1}] \big| \leq \epsilon_i \, \ecty(f_i|\mcal{F}_{i-1}) \, \ecty(g|\mcal{F}_{i-1}) ,\]
which yields upon integrating:
\begin{multline}\label{cal3329}
|\EE[f_ig]| \leq \epsilon_i \int \ecty(f_i|\mcal{F}_{i-1}) \ecty(g|\mcal{F}_{i-1}) \,\dx\Pr \\
\footrel{\text{CS}}{\leq} \epsilon_i \sqrt{\int \Var(f_i|\mcal{F}_{i-1}) \, \dx\Pr} \sqrt{\int \Var(g|\mcal{F}_{i-1}) \, \dx\Pr} \\
= \epsilon_i \sqrt{V_i} \sqrt{\Var(g) - \Var(g^{\mcal{F}_{i-1}})} \leq \epsilon_i \sqrt{V_i} \ecty(g) .
\end{multline}

For the second term, under~$\Pr[\Bcdot|\ab \mcal{F}_i]$,
$g$ only depends on~$Y$, and $\EE[\tilde{f}_i|Y]
\equiv \tilde{f}_i - \tilde{f}^*_i$ as we noticed just before Claim~\ref{clm1908},
so $\EE[\tilde{f}_ig|\mcal{F}_i] = \EE[{(\tilde{f}_i - \tilde{f}^*_i)}g | \mcal{F}_i]$,
which yields upon integrating:
\[\label{for3330} |\EE[\tilde{f}_ig]|
= \big| \EE \big[ (\tilde{f}_i - \tilde{f}^*_i) g \big] \big|
\footrel{\text{CS}}{\leq} \ecty(\tilde{f}_i - \tilde{f}^*_i) \ecty(g) = \sqrt{\tilde{V}_i - \tilde{V}^*_i} \ecty(g) ,\]
the last equality coming from Claim~\ref{clm1908}.
Then it just remains to combine~(\ref{cal3329}) and~(\ref{for3330})
to get~(\ref{for3331}).\linebreak[1]\strut
\end{proof}\endgroup

\begingroup\def\proofname{Proof of Theorem~\ref{thm6413}}\begin{proof}
First, thanks to a by now classical approximation argument
we may assume that $I = \{1,\ldots,N\}$
and $J = \{1,\ldots,M\}$. Denote $\mcal{F} \coloneqq \sigma(\vec{X}_I)$,
resp.\ $\mcal{G} \coloneqq \sigma(\vec{Y}_J)$; our goal is to prove
that for all~$f\in\ldb(\mcal{F})$, all~$g\in\ldb(\mcal{G})$,
one has $|\EE[fg]| \leq ( \VERT \boldepsilon \VERT \wedge 1 ) \ecty(f) \ecty(g)$.
We will use the same trick as in our alternative proof of Theorem~\ref{thm5750}:
by the definition of conditional expectation and the Cauchy--Schwarz inequality,
proving the inequality above is equivalent to showing that for all~$f\in\ldb(\mcal{F})$,
\[ \Var\big( f^{\mcal{G}} \big) \leq \big( \VERT \boldepsilon \VERT^2 \wedge 1 \big) \Var(f) ,\]
which, by associativity of variance, is in turn equivalent to:
\[\label{f3451}
\Var\big( f - f^{\mcal{G}} \big) \geq \big( 1-\VERT\boldepsilon\VERT^2 \big)_+ \Var(f) .\]

For~$0\leq i\leq N$, resp.~$0\leq j\leq M$, define
$\mcal{F}_i = \sigma(X_1,\ldots,X_i)$, resp. $\mcal{G}_j = \sigma(Y_1,\ldots,Y_j)$.
For all~$0 \leq j \leq M$, define
\[ f^j \coloneqq f - \EE[f|\mcal{G}_j] ,\]
and for all~$1 \leq i \leq N$, define moreover
\[ f^j_i \coloneqq f^{\mcal{G}_j\vee\mcal{F}_i} - \EE[f|\mcal{G}_j\vee\mcal{F}_{i-1}] .\]
Denote $V^j \coloneqq \Var(f^j)$, resp.\ $V^j_i \coloneqq \Var(f^j_i)$.
For fixed~$j$, the~$f^j_i$ are pairwise orthogonal
(again by the argument in the proof of Proposition~\ref{pro5527})
and their sum is equal to~$f^j$, so:
\[\label{for3043} V^j = \sum_{i=1}^N V^j_i .\]
Thus, with this notation our goal~(\ref{f3451}) becomes:
\[\label{for7254}
\sum_{i=1}^N V^M_i \geq \big( 1-\VERT\boldepsilon\VERT^2 \big)_{\!+\,} \sum_{i=1}^N V^0_i .\]

The main tool to prove~(\ref{for7254}) will be Lemma~\ref{lem7265}.
Actually the rough formula~(\ref{for1302}) is quite impratical,
so we introduce a \emph{linearized} version of it:
for each~$1\leq i\leq N$ take some $\alpha_i > 0$
(which for the time being is arbitrary), then by the Cauchy--Schwarz inequality,
(\ref{for1302}) implies that:
\[\label{for3261}
V^*_i \geq (1-\epsilon_i^2) V_i - \frac{\epsilon_i V_i}{\alpha_i} \sum_{i'>i} \epsilon_{i'}\alpha_{i'}
- \epsilon_i\alpha_i \sum_{i'>i} \frac{\epsilon_{i'}V_{i'}}{\alpha_{i'}} .\]

\begin{Rmk}
(\ref{for3261}) is devised so that its right-hand side is exactly the same
as in~(\ref{for1302}) if $V_i \propto \alpha_i^2\enskip\forall i$.
\end{Rmk}

Let us reason conditionally to~$\mcal{G}_{j-1}$ for a few lines.
Under this conditioning, call~$\dot{\mcal{F}}_i \coloneqq \sigma(X_1,\ldots,X_i)$,
resp.\ $\dot{\mcal{F}}^*_i \coloneqq \sigma(X_1,\ldots,X_i,Y_j)$, and $\dot{f} \coloneqq f^{j-1}$.
Then $\dot{f}$ is an $\ldb(\vec{X})$ function,
so we are in situation of applying Lemma~\ref{lem7265} to the functions
\begin{eqnarray} \dot{f}_i &\coloneqq& \dot{f}^{\dot{\mcal{F}}_i} - \EE\big[\dot{f}\big|\dot{\mcal{F}}_{i-1}\big] \\
\text{\llap{and \qquad}{}} \dot{f}^*_i &\coloneqq& \dot{f}^{\dot{\mcal{F}}^*_i} - \EE\big[\dot{f}\big|\dot{\mcal{F}}^*_{i-1}\big] .
\end{eqnarray}
%
But in fact we already know these functions:
namely, $\dot{f}_i = f^{j-1}_i$ and $\dot{f}^*_i = f^j_i$.
Then, applying the linearized version~(\ref{for3261}) of Lemma~\ref{lem7265}
\[ \Var \big( f^j_i|\mcal{G}_{j-1} \big) \geq \\
\bigg( 1-\epsilon_{ij}^2 - \frac{\epsilon_{ij}}{\alpha_i} \sum_{i'>i} \epsilon_{i'j}\alpha_{i'} \bigg)
\Var \big( f^{j-1}_i|\mcal{G}_{j-1} \big)
- \epsilon_{ij}\alpha_i \sum_{i'>i} \frac{\epsilon_{i'j}\Var\big( f^{j-1}_{i'}|\mcal{G}_{j-1} \big)}{\alpha_{i'}} , \]
whence upon integrating:
\[\label{for4382}
V^j_i \geq (1-\epsilon_{ij}^2) V^{j-1}_i
- \Big(\sum_{i'>i} \epsilon_{i'j}\alpha_{i'}\Big) \frac{\epsilon_{ij} V^{j-1}_i}{\alpha_i}
- \epsilon_{ij}\alpha_i \sum_{i'>i} \frac{\epsilon_{i'j}V^{j-1}_{i'}}{\alpha_{i'}} .\]

By Equation~(\ref{for4382}), we have transformed our initial problem
into a purely abstract operator problem, \emph{posed in an $L^1$ setting}.
%
To handle it, we need a little notation.
Call~$L^1(I)$ the set of real functions on~$I$,
endowed with the $L^1$ norm
\[ \big\| (v_i)_{i\in I} \big\|_1 \coloneqq \sum_{i\in I} |v_i| .\]
The dual space of~$L^1(I)$ is made of the linear forms
$\ell \colon (v_i)_{i\in I} \mapsto \sum \ell_i v_i$,
equipped with the $L^{\infty}$ norm
\[ \| \ell \|_{\infty} \coloneqq \sup_{i\in I} |\ell_i| .\]
We shall write ``$L^1(I) \ni v \geq 0$'' to mean that all the entries of~$v$ are nonnegative,
and ``$(L^1(I))' \ni \ell \geq 0$'' to mean that
$(v \geq 0) \Rightarrow (\ell v \geq 0)$, which is equivalent to say that
all the~$\ell_i$ are nonnegative.
%
Now I claim the following lemma, whose proof is postponed:
\begin{Lem}\label{lem5453}
Suppose given some nonnegative numbers $V^j_i$ for
$(i,j) \in \{1,\ldots,N\} \times \{0,\ldots,M\}$,
such that Equation~(\ref{for4382}) is satisfied for all~$i,j$.
Call~$\mcal{L}$ the nonnegative linear form on~$L^1(I)$ defined by
\[ \mcal{L}v = \sum_{\substack{j\in J\\i,i'\in I}} \frac{\alpha_{i'}}{\alpha_i} \epsilon_{ij}\epsilon_{i'j} v_i ,\]
and assume $\|\mcal{L}\|_{\infty} \leq 1$,
then:
\[\label{for6376}
\sum_{i=1}^N V^M_i \geq \sum_{i=1}^N V^0_i - \mcal{L}\big( (V^0_i)_{i\in I} \big) .\]
\end{Lem}

Lemma~\ref{lem5453} has the following immediate
\begin{Cor}\label{cor1307}
Suppose given some nonnegative numbers $V^j_i$ for
$(i,j) \in \{1,\ldots,N\} \times \{0,\ldots,M\}$,
such that Equation~(\ref{for4382}) is satisfied for all~$i,j$,
then:
\[\label{for1338} \sum_{i=1}^N V^M_i \geq
\bigg( 1 - \sup_{i\in I} \sum_{\substack{j\in J\\i'\in I}} \frac{\alpha_{i'}}{\alpha_i} \epsilon_{ij}\epsilon_{i'j}
\bigg)_{\!\!+\,\,} \sum_{i=1}^N V^0_i .\]
\end{Cor}

Now we finish the proof of Theorem~\ref{thm6413}:
thanks to Corollary~\ref{cor1307}
we have proved that~(\ref{for1338}) stands true in our situation
for \emph{any} choice of positive $(\alpha_i)_{i\in I}$.
The last step then consists in optimizing that choice.
Denote ``$\alpha>0$'' to mean that all the~$\alpha_i$ are positive. One has:
\begin{multline}
\inf_{\alpha>0} \sup_{i\in I} \sum_{\substack{j\in J\\i'\in I}} \frac{\alpha_{i'}}{\alpha_i} \epsilon_{ij}\epsilon_{i'j}
= \inf \big\{ \lambda \geq 0 \colon
(\exists \alpha > 0) (\forall i) \big( \sum_{\substack{j\in J\\i'\in I}} \epsilon_{ij}\epsilon_{i'j}\alpha_{i'}
\leq \lambda \alpha_i \big) \big\} \\
= \inf \big\{ \lambda \geq 0 \colon (\exists \alpha > 0)
(\boldepsilon\boldepsilon^* \alpha \leq \lambda\alpha) \big\}.
\end{multline}
But $\boldepsilon\boldepsilon^*$ is a nonnegative operator on~$L^2(I)$
(I mean, when seen as a matrix all its entries are nonnegative),
so by Lemma~\ref{l3542} in appendix:
%
\[\label{for2029}
\inf \big\{ \lambda\geq 0 \colon (\exists \alpha > 0) (\boldepsilon\boldepsilon^* \alpha \leq \lambda\alpha) \big\}
= \rho(\boldepsilon\boldepsilon^*) = \VERT \boldepsilon \VERT^2 .\]
This ends the proof of Theorem~\ref{thm6413}.
\end{proof}\endgroup

\begingroup\def\proofname{Proof of Lemma~\ref{lem5453}}\begin{proof}
We prove Lemma~\ref{lem5453} by induction on~$M$.
The case $M=0$ is trivial.
Suppose $M\geq1$ and assume the result is true for~${(M-1)}$.
We generalize the notation $\mcal{L}$
by defining, for~$\bullet \in \{\text{\textvisiblespace},1,*\}$,
\[ \mcal{L}^{\bullet}v = \sum_{\substack{j\in J^{\bullet}\\i,i'\in I}}
\frac{\alpha_{i'}}{\alpha_i} \epsilon_{ij}\epsilon_{i'j} v_i ,\]
with $J^1 = \{1\}$, resp.\ $J^* = \{2,\ldots,M\}$,
so that $\mcal{L} = \mcal{L}^1 + \mcal{L}^*$.
Notice that $\|\mcal{L}^*\|_{\infty}\leq 1$ since ${\|\mcal{L}\|_{\infty}\leq1}$.
For all~$i \in I$, define
\[ \check{V}^1_i = (1-\epsilon_{i1}^2) V^0_i - \frac{\epsilon_{i1} V^0_i}{\alpha_i} \sum_{i'>i} \epsilon_{i'1}\alpha_{i'}
- \epsilon_{i1}\alpha_i \sum_{i'>i} \frac{\epsilon_{i'1}V^0_{i'}}{\alpha_{i'}} ,\]
which is the value that $V^1_i$ would take if there were equality in~(\ref{for4382})
for $j=1$.
%
With that notation, (\ref{for4382}) writes
\[\label{for6134} \big( V^1_i - \check{V}^1_i \big)_{i\in I} \geq 0 ,\]
and by induction hypothesis we have:
\[\label{for6135} \sum_{i=1}^N V^M_i \geq \sum_{i=1}^N V^1_i - \mcal{L}^* \big( (V^1_i)_{i\in I} \big) .\]

Introducing the~$\check{V}^1_i$, we have therefore the following chain of inequalities:
\begin{multline}
\sum_{i=1}^N V^M_i \footrel{(\ref{for6135})}{\geq}
\sum_{i=1}^N V^1_i - \mcal{L}^* \big( (V^1_i)_{i\in I} \big) \\
%
\footrel{(\ref{for6134})}{=} \sum_{i=1}^N \check{V}^1_i + \big\| \big( V^1_i - \check{V}^1_i \big)_{i\in I} \big\|_1
- \mcal{L}^* \big( (\check{V}^1_i)_{i\in I} \big) - \mcal{L}^* \big( (V^1_i - \check{V}^1_i)_{i\in I} \big) \\
%
\footrel{\|\mcal{L}^*\|_\infty\leq1}{\geq} \sum_{i=1}^N \check{V}^1_i - \mcal{L}^* \big( (\check{V}^1_i)_{i\in I} \big)
%
= \sum_{i=1}^N V^0_i - \mcal{L}^1 \big( (V^0_i)_{i\in I} \big)
- \mcal{L}^* \big( (\check{V}^1_i)_{i\in I} \big) \\
%
\footrel{\substack{\check{V}^1\leq V^0\\\mcal{L}^*\geq0}}{\geq}
\sum_{i=1}^N V^0_i - \mcal{L}^1 \big( (V^0_i)_{i\in I} \big)
- \mcal{L}^* \big( (V^0_i)_{i\in I} \big)
%
= \sum_{i=1}^N V^0_i - \mcal{L} \big( (V^0_i)_{i\in I} \big),
\end{multline}
%
so~(\ref{for6376}) is true for~$M$, whence the lemma by induction.
\end{proof}\endgroup

\begin{Rmk}
Our proof of Theorem~\ref{thm6413}
handled the~$X_i$ and the~$Y_j$ in a fully nonsymmetric way,
since we began with putting orders on~$I$ and~$J$,
which orders played a crucial role in the decomposition of~$f$.
%
Yet the bound~(\ref{for3009}) obtained \emph{is} obviously symmetric
by re-labelling the basic variables%
—and this is not due to having proceeded to any `re-symmetrization' step\dots\ 
To date I have no simple explanation for this `coincidence'.
\end{Rmk}


\subsection{`\texorpdfstring{$\ZZ$}{Z}~against~\texorpdfstring{$\ZZ$}{Z}' tensorization}

The proof of the `$N$~against~$M$' theorem was quite more technical than
that of the `$N$~against~$1$' theorem;
because of that, in order to get tractable computations
we had to use suboptimal inequalities at two places:
%
\begin{itemize}
\item Claim~\ref{clm1959} is suboptimal: it has indeed the same shortcoming
as Proposition~\ref{pro5527} exhibited compared to Theorem~\ref{thm5750},
namely, it does not `recycle the losses' occurring
when one makes $g$ covariate with both~$f_i$ and~$\tilde{f}_i$
(cf.\ the discussion on page~\pageref{lbl91},
just after the proof of Proposition~\ref{pro5527}).
%
\item Our linearization technique is suboptimal in general,
even after optimizing the~$\alpha_i$.
In fact, as we said before, Inequality~(\ref{for3261}) is optimal
if and only if one has $V_i \propto \alpha_i$;
thus, for~(\ref{for4382}) to be always optimal,
one has to have $V^j_i \propto \alpha_i$
for \emph{all} $j$, with the \emph{same} values for the~$\alpha_i$.
This would imply that all the sequences $(V_i^j)_{0\leq i<n}$
are proportional, which is not true in general.
\end{itemize}

So, Theorem~\ref{thm6413} is certainly not optimal%
\footnote{Though, as we will see in \S~\ref{parAsymptoticOptimality},
it is `asymptotically optimal'.}%
—this is confirmed by the example of \S~\ref{parIllustrationOfTheProofs}.
Nonetheless, there is one particular case in which
an alternative reasoning yields an optimal bound%
\footnote{The bound's being optimal shall be proved by Theorem~\ref{t6657}.}.
This case is when some symmetries in the decorrelation hypotheses
allow us to transform the original two-parameter problem (indexed by~$I\times J$)
into a one-parameter problem (indexed by~$\ZZ$).
Let us state and prove the corresponding result:
%
\begin{Thm}[`$\ZZ$~against~$\ZZ$' theorem]\label{thm0119}
Let~$I$ and~$J$ be sets isomorphic to~$\ZZ$,
and let~$(X_i)_{i\in I}$ and~$(Y_j)_{j\in J}$ be random variables
such that, $\mcal{M}$ denoting the $\sigma$-metalgebra they generate,
one has for all~$i,j\in\ZZ$
\[ \{X_i : Y_j\}_{\mcal{M}} \leq \epsilon(j-i) \]
for some function~$\epsilon \colon \ZZ \to [0,1]$.

Then
\[ \big\{ \vec{X}_I : \vec{Y}_J \big\} \leq \bar{\epsilon} ,\]
where $\bar{\epsilon} \in [0,1]$ is characterized by:
\[\label{f6444} \arcsin \bar{\epsilon} =
\Big( \sum_{z\in\ZZ} \arcsin \epsilon(z) \Big) \wedge \frac{\pi}{2} .\]
\end{Thm}

\begin{Rmk}
If we apply Theorem~\ref{thm6413} to the situation above, we find
${\{\vec{X}_I:\vec{Y}_J\}} \ab \leq \big( \sum_{z\in\ZZ}\epsilon(z) \big) \wedge 1$
(cf.\ \S~\ref{parPractical}).
The latter expression is always $\geq \bar{\epsilon}$
because of the concavity of the function~$\sin(\Bcdot\wedge\frac{\pi}{2})$ on~$\RR_+$,
and even $>\bar\epsilon$ if $\bar{\epsilon}\neq0,1$;
so, when it is applicable, Theorem~\ref{thm0119} is strictly stronger than Theorem~\ref{thm6413}.
\end{Rmk}

\begin{proof}
Let~$f$ and~$g$ be resp.\ $\vec{X}_I$- and $\vec{Y}_J$-measurable $\ldb$ functions.
Denote $\mcal{F} \coloneqq \sigma(\vec{X})$, resp.\ $\mcal{G} \coloneqq \sigma(\vec{Y})$,
and for~$i\in\ZZ$, resp.~$j\in\ZZ$,
denote $\mcal{F}_i \coloneqq \bigvee_{i'\leq i} \sigma(X_{i'})$,
resp.\ $\mcal{G}_j \coloneqq \bigvee_{j'\leq j} \sigma(Y_{j'})$.
For~$(i,j)\in\ZZ\times\ZZ$, define
\[\label{f9113}
f_i^j \coloneqq f^{\mcal{G}_j\vee\mcal{F}_i} - \EE[f|\mcal{G}_j\vee\mcal{F}_{i-1}] \]
and\footnote{Beware: the definition of~$g_j^i$
is \emph{not} analogous to the definition of~$f_i^j$\,!}
\[ g_j^i \coloneqq g^{\mcal{G}_j} - \EE[g^{\mcal{G}_j}|\mcal{G}_{j-1}\vee\mcal{F}_i] .\]
Denote $V\coloneqq\Var(f)$, $W\coloneqq\Var(g)$, $V_i^j\coloneqq\Var(f_i^j)$, $W_j^i\coloneqq\Var(g_j^i)$;
also denote
\[ S_{ij} \coloneqq \EE[f_i^{j-1}g_j^{i-1}] .\]

Our auxiliary functions were devised so that
\begin{Clm}\label{clm4070}
Provided the sum in the right-hand side is absolutely convergent,
\[\label{for5163} \EE[fg] = \sum_{i,j} S_{ij} .\]
\end{Clm}

\begingroup\def\proofname{Proof of Claim~\ref{clm4070}}\begin{proof}
First define $\bar{f} \coloneqq f^{\mcal{G}}$,
so that $\bar{f}$ is $\mcal{G}$-measurable and
$\EE[fg] = \EE[\bar{f}g]$.
For~$j \ab {\in \ZZ}$, define
$g_j \coloneqq g^{\mcal{G}_j} - \EE[g|\mcal{G}_{j-1}]$,
resp.\ $\bar{f}_j \coloneqq {\bar f}^{\mcal{G}_j} - \EE[\bar{f}|\mcal{G}_{j-1}]$:
we have $g=\sum_jg_j$ and $\bar{f}=\sum_j\bar{f}_j$,
which are the respective decompositions of~$g$ and~$\bar{f}$
on the same basis of orthogonal subspaces of~$\ldb(\mcal{G})$,
so $\EE[fg] = \sum_{j} \EE[\bar{f}_jg_j]$.
The terms of the right-hand side of that formula are unchanged upon replacing
$\bar{f}_j$ by~$f^{j-1} \coloneqq f - \EE[f|\mcal{G}_{j-1}]$,
since $\EE[(f^{j-1} - \bar{f}_j)g_j]$ is zero%
—the function~$(f^{j-1} - \bar{f}_j)$ is indeed equal to~$(f-\EE[f|\mcal{G}_j])$,
which is centered conditionally to~$\mcal{G}_j$,
while $g_j$ is $\mcal{G}_j$-measurable.
In the end we have:
\[\label{for5102} \EE[fg] = \sum_j \EE[f^{j-1}g_j] .\]

So in a first step we have decomposed $\EE[fg]$
into a sum indexed by~$j$. Now we decompose each term of that sum
into a sum indexed by~$i$.
Let us reason conditionally to~$\mcal{G}_{j-1}$.
Then $f^{j-1}$ is an $\ldb(\mcal{F})$ function
and $g_j$ is in~$\ldb(Y_j)$.
%
We compute $\EE[f^{j-1}g_j]$ as in the first step of this proof:
first we replace $g_j$ by~$\bar{g}_j \coloneqq (g_j)^{\mcal{F}}$;
then we decompose $f^{j-1} = \sum_i f^{j-1}_i$ and $\bar{g}_j = \sum_i \bar{g}_{ji}$,
with $f^{j-1}_i \coloneqq (f^{j-1})^{\mcal{F}_i} - \EE[f^{j-1}|\mcal{F}_{i-1}]$%
%
\footnote{Notation is consistent: this $f^{j-1}_i$ is indeed the same
as the~$f^{j-1}_i$ defined by~(\ref{f9113}),
since we are reasoning conditionally to~$\mcal{G}_{j-1}$.},
%
resp.\ $\bar{g}_{ji} \coloneqq \bar{g}_j^{\mcal{F}_i} - \EE[\bar{g}_j|\mcal{F}_{i-1}]$,
and by orthogonal decomposition
we get $\EE[f^{j-1}g_j] = \sum_i \EE[f^{j-1}_i\bar{g}_{ji}]$;
then we conclude by saying that $\EE[f^{j-1}_i\bar{g}_{ji}]$
is actually equal to~$\EE[f^{j-1}_ig^{i-1}_j]$,
since $(g^{i-1}_j - \bar{g}_{ji})$ is centered conditionally to~$\mcal{F}_i$
while $f^{j-1}_i$ is $\mcal{F}_i$-measurable.
In the end we have obtained
\[\label{for5103} \EE[f^{j-1}g_j] = \sum_i \EE[f^{j-1}_ig^{i-1}_j] ,\]
which combined with~(\ref{for5102}) yields~(\ref{for5163}).
\end{proof}\endgroup

So we have expressed $\EE[fg]$ as a function of the~$S_{ij}$.
It is also possible to `read' the values of~$V$ and~$W$ from the~$V_i^j$,
resp.\ from the~$W_j^i$, \emph{via} the formulas:
\begin{eqnarray}
\label{for0923a} V &=& \overset{\nearrow}{\lim_{j\longto-\infty}} \Big( \sum_i V_i^j \Big) ;\\
\label{for0923b} W &=& \sum_j \Big( \overset{\nearrow}{\lim_{i\longto-\infty}} W_j^i \Big) .
\end{eqnarray}

Now we are looking for relations between the~$V_i^j$, the~$W_j^i$ and the~$S_{ij}$.
The first relation comes from the decorrelation hypothesis:
conditionally to~$\mcal{G}_{j-1} \vee \mcal{F}_{i-1}$,
$f_i^{j-1}$ is in~$\ldb(X_i)$, resp.~$g_j^{i-1}$ is in~$\ldb(Y_j)$,
and ${\{X_i:Y_j\}} \leq \epsilon(j-i)$, so:
\[\label{for9472a} |S_{ij}| \leq \epsilon(j-i) \sqrt{V_i^{j-1}W_j^{i-1}} .\]

The second relation means that a large value of~$|S_{ij}|$ forces $W_j^i$ to diminish.
To state it, we observe that, since $f_i^{j-1}$ is $(\mcal{G}_{j-1}\vee\mcal{F}_i)$-measurable,
$S_{ij} = \EE[f_i^{j-1}(g_j^{i-1})^{\mcal{G}_{j-1}\vee\mcal{F}_i}]$,
so by the Cauchy--Schwarz inequality
$|S_{ij}| \leq \ecty(f_i^{j-1}) \*
\ecty\big( (g_j^{i-1})^{\mcal{G}_{j-1}\vee\mcal{F}_i} \big)$.
Moreover, since $g_j^{i-1} - (g_j^{i-1})^{\mcal{G}_{j-1}\vee\mcal{F}_i} = g_j^i$,
one has by orthogonality $\Var\big( (g_j^{i-1})^{\mcal{G}_{j-1}\vee\mcal{F}_i} \big) =
\Var(g_j^{i-1}) - \Var(g_j^i)$,
so our inequality becomes
\[\label{for9472b} |S_{ij}| \leq \sqrt{V_i^{j-1}}\sqrt{W_j^{i-1}-W_j^i} \]
(where it is understood that $W_j^i \leq W_j^{i-1}$),
or more eloquently
\[\label{for9472b'} W^i_j \leq W^{i-1}_j - (S_{ij})^2/V^{j-1}_i \]
provided $V_i^{j-1} > 0$.

The third and last relation means, on the other hand,
that a large value of~$\big|\sum_{i'>i} S_{i'j}\big|$ forces~$\sum_{i'>i} V_i^j$ to diminish.
To state it, we denote
\[ \tilde{f}^j_i \coloneqq f - f^{\mcal{G}_j\vee\mcal{F}_i}
= \sum_{i'>i} f^j_{i'} ,\]
whose variance is $\Var(\tilde{f}^j_i) = \sum_{i'>i} \Var(f_{i'}^j)$
since the~$f_{i'}^j$ are pairwise orthogonal.
One has
\[ \sum_{i'>i} S_{i'j} = \sum_{i'>i} \EE[f_{i'}^{j-1} g_j^{i'-1}]
= \sum_{i'>i} \EE [ f_{i'}^{j-1} g_j^i ] = \EE[\tilde{f}_i^{j-1} g_j^i]
= \EE[(\tilde{f}_i^{j-1})^{\mcal{G}_j\vee\mcal{F}_i} g_j^i] , \]
so by the Cauchy--Schwarz inequality,
\[ \Big| \sum_{i'>i} S_{i'j} \Big| \leq
\ecty\big((\tilde{f}^{j-1}_i)^{\mcal{G}_j\vee\mcal{F}_i}\big) \ecty(g_j^i) .\]
Since $\tilde{f}^{j-1}_i - (\tilde{f}^{j-1}_i)^{\mcal{G}_j\vee\mcal{F}_i} = \tilde{f}^j_i$,
one has by orthogonality
\[ \Var\big((\tilde{f}^{j-1}_i)^{\mcal{G}_j\vee\mcal{F}_i}\big) =
\Var(\tilde{f}^{j-1}_i) - \Var(\tilde{f}^j_i) ,\]
so our inequality becomes
\[\label{for9472c}
\Big| \sum_{i'>i} S_{i'j} \Big| \leq
\sqrt{\sum_{i'>i} V_{i'}^{j-1} - \sum_{i'>i} V_{i'}^j} \sqrt{W_j^i} ,\]
or more eloquently:
\[\label{for9472c'} \sum_{i'>i} V_{i'}^j \leq
\sum_{i'>i} V_{i'}^{j-1} - \Big( \sum_{i'>i} S_{i'j} \Big)^{\!2} \mathbin{\Big/} W_j^i .\]

So, we have transformed our initial probabilistic problem into the following analytic one:
let~$\mbf{A}$ be an array indexed by~$\ZZ\times\ZZ$,
each entry $(i,j)$ of which contains three numbers~$V_i^j\geq 0$, $W_j^i\geq 0$ and~$S_{ij}$,
satisfying~(\ref{for9472a}), (\ref{for9472b}) and~(\ref{for9472c})%
—we will say such an array is \emph{correct}.
We define $V$ by~(\ref{for0923a}) and~$W$ by~(\ref{for0923b}),
and we set $S = \sum_{i,j} S_{ij}$ (provided it makes sense);
our goal is to get a bound of the form
``$|S| \leq \bar{\epsilon}\sqrt{VW}$'',
with~$\bar{\epsilon}$ only depending on~$\epsilon(\Bcdot)$.

Note that \emph{A priori} some problems of summability can arise from $\mbf{A}$'s being infinite,
for instance to check~(\ref{for9472c'}) or to define $S$.
However, in the situations which are of interest to us,
we can restrict to cases in which $\mbf{A}$ is of nice particular form.
To do this, we first approximate $f$ in~$\ldb(\vec{X})$, resp.~$g$ in~$\ldb(\vec{Y})$,
by a function depending only on a finite number of~$X_i$, resp.\ of~$Y_j$
—say, we assume $f$ is $\vec{X}_{\dot{I}}$-measurable
and $g$ is $\vec{Y}_{\dot{J}}$-measurable for finite $\dot{I} \subset I, \dot{J} \subset J$.
%
Then, we define a new model $(\tilde{X}_i)_{i\in\ZZ}, (\tilde{Y}_j)_{j\in\ZZ}$
by~$\tilde{X}_i = X_i$ for $i\in\dot{I}$, resp.\ $\tilde{Y}_j = Y_j$ for $j\in\dot{J}$,
and $\tilde{X}_i, \tilde{Y}_j = \partial$ for $i\notin\dot{I}, j\notin\dot{J}$,
$\partial$ being some cemetery point.
This new model still gives a correct array,
for which $S/\sqrt{VW}$ is arbitrarily close to the initial value of
$\EE[fg]/\ecty(f)\ecty(g)$;
and the new array is of the following form, which we will call \emph{compact},
for which all the quantities of interest are well defined:
\begin{itemize}
\item $V^j_i$ is zero as soon as $i\notin\dot{I}$, and it does not depend on~$j$
for $j < \min\dot{J}$, nor for $j \geq \max\dot{J}$;
\item Similarly, $W^i_j$ is zero as soon as $j\notin\dot{J}$, and it does not depend on~$i$
for $i < \min\dot{I}$, nor for $i \geq \max\dot{I}$;
\item $S_{ij}$ is zero as soon as $(i,j)\notin\dot{I}\times\dot{J}$.
(This condition automatically follows from the first two if the array is correct).
\end{itemize}

We define the following operations on arrays:
\begin{Def}\strut\begin{itemize}
\item For~$z\in\ZZ$, we define the \emph{translation operator} $\tau^z$ on arrays
such that, if the entries of~$\mbf{A}$ at~$(i,j)$ are $V_i^j,W_j^i,S_{ij}$,
the entries of~$\tau^z\mbf{A}$ at~$(i,j)$ are $V_{i+z}^{j+z},W_{j+z}^{i+z},S_{(i+z)(j+z)}$.
\item For~$\grave{\mbf{A}}$ and~$\acute{\mbf{A}}$ two arrays
with entries $\grave{V}_i^j, \ab \grave{W}_j^i, \ab \grave{S}_{ij}$,
resp.~$\acute{V}_i^j, \ab \text{etc.}$,
for~$\alpha, \beta$ two real numbers,
we define the linear combination $\alpha\grave{\mbf{A}} + \beta\acute{\mbf{A}}$
as the array with entries $\alpha\grave{V}_i^j+\beta\acute{V}_i^j,
\alpha\grave{W}_j^i+\beta\acute{W}_j^i,\text{etc.}$.
\end{itemize}\end{Def}

\begin{Lem}\label{lem0513}
Correct arrays are stable by translations and by nonnegative linear combinations,
i.e., if $\mbf{A}$ and~$\mbf{B}$ are correct arrays,
then for all~$z\in\ZZ$ and~$\alpha,\beta \geq 0$,
$\tau^z\mbf{A}$ and~$\alpha\mbf{A} + \beta\mbf{B}$
are correct too.
\end{Lem}

\begingroup\def\proofname{Proof of Lemma~\ref{lem0513}}\begin{proof}
Recall that being correct means satisfying~(\ref{for9472a}), (\ref{for9472b}) and~(\ref{for9472c}).
These conditions are trivially stable by multiplication by a nonnegative constant
and by translation%
\footnote{Getting stability of Condition~(\ref{for9472a}) by translation
is actually the only place where the symmetries of the problem are used.}.
It remains to see that they are stable by addition.
The technique being the same for all three inequalities,
we just treat the case of~(\ref{for9472b}).
%
Stability of this condition by addition is a consequence of the following inequality
(which is in fact a particuliar case of the \emph{Brunn--Minkowski inequality},
see~\cite{B-M_inequality}):
%
\begin{Lem}\label{lem0514} For all~$a_1, b_1, a_2, b_2 \geq 0$,
\[\label{f1709} \sqrt{(a_1+a_2)(b_1+b_2)} \geq \sqrt{a_1b_1} + \sqrt{a_2b_2} .\]
\end{Lem}
%
\begingroup\def\proofname{Proof of Lemma~\ref{lem0514}}\begin{proof}
Take squares on both sides of~(\ref{f1709})
and notice that ${(a_1+a_2)}\ab {(b_1+b_2)} \ab - {\big(\sqrt{a_1b_1} + \sqrt{a_2b_2}\big)^2}
\ab = a_1b_2 + a_2b_1 - 2\sqrt{a_1b_1a_2b_2} \ab =
{\big(\sqrt{a_1b_2} - \sqrt{a_2b_1}\big)^2} \ab \geq 0$.
\end{proof}\endgroup

For~$\grave{\mbf{A}}$ and~$\acute{\mbf{A}}$ two correct arrays satisfying~(\ref{for9472b}),
applying~(\ref{f1709}) with ${a_1 = \acute{V}^{j-1}_i}, a_2 = \grave{V}^{j-1}_i,
b_1 = \acute{W}^{i-1}_j - \acute{W}^i_j, b_2 = \grave{W}^{i-1}_j - \grave{W}^i_j$,
we get:
\begin{multline} \big|\grave{S}_{ij}+\acute{S}_{ij}\big| \leq
|\grave{S}_{ij}| + |\acute{S}_{ij}| \leq
\sqrt{\grave{V}_i^{j-1}}\sqrt{\grave{W}_j^{i-1}-\grave{W}_j^i}
+ \sqrt{\acute{V}_i^{j-1}}\sqrt{\acute{W}_j^{i-1}-\acute{W}_j^i} \\
\leq \sqrt{\grave{V}_i^{j-1}+\acute{V}_i^{j-1}}
\sqrt{(\grave{W}_j^{i-1}+\acute{W}_j^{i-1}) - (\grave{W}_j^i+\acute{W}_j^i)} ,
\end{multline}
so~(\ref{for9472b}) is still valid for~$(\grave{\mbf{A}}+\acute{\mbf{A}})$.
\end{proof}\endgroup

Now, thanks to Lemma~\ref{lem0513} we will reduce our problem on~$(\ZZ\times\ZZ)$-arrays
into a problem on~$\ZZ$-arrays.
Suppose $\mbf{A}$ is a correct array with certain values of~$V$, $W$ and~$S$.
Then, for~$k \geq 0$, the array
\[\label{for3010} \mbf{A}_k = \frac{1}{2k+1} \sum_{z=-k}^{k} \tau^z\mbf{A} \]
is correct too, with the same values of~$V$, $W$ and~$S$ as~$\mbf{A}$.
Now when~${k \longto \infty}$, $\mbf{A}_k$ `looks more and more like a Toeplitz array',
that is, an array whose entries at~$(i,j)$ only depend on~$(j-i)$.
To state it rigorously, we need some definitions:
%
\begin{Def}\label{def3111}\strut
\begin{itemize}
\item Here, a \emph{Toeplitz array} will mean a $\ZZ\times\ZZ$ array whose entries at~$(i,j)$
only depend on~$(j-i)$.
For such an array, for~$z\in\ZZ$ we denote by~$V_{(z)}, W_{(z)}, S_{(z)}$ the quantities
characterized by~$V_i^j = V_{(j-i)}$, etc..
%
\item Actually we can always assume our Toeplitz array is \emph{Toeplitz compact},
which means that there exists some $z^-\leq z^+$ such that:
\begin{itemize}
\item $V_{(z)}$ does not depend on~$z$ for $z<z^-$, nor for $z\geq z^+$;
\item $W_{(z)}$ does not depend on~$z$ for $z\leq z^-$, nor for $z>z^+$;
\item $S_{(z)}$ is zero as soon as $z<z^-$ or $z>z^+$.
\end{itemize}
%
\item For a compact Toeplitz array, we define $v, w, s$
as `renormalized versions' of~$V, W, S$:
\begin{eqnarray}
\label{for3162a} v &\coloneqq& V_{(z<z^-)} ;\\
\label{for3162b} w &\coloneqq& W_{(z>z^+)} ;\\
\label{for3162c} s &\coloneqq& \sum_{z\in\ZZ} S_{(z)} .
\end{eqnarray}
%
\item A Toeplitz array is said to be \emph{correct} if it is correct when seen as an ordinary array.
For a Toeplitz array, Equations~(\ref{for9472a}), (\ref{for9472b'}) and~(\ref{for9472c'})
become respectively%
\footnote{Note that the way (\ref{for3534}) follows from~(\ref{for9472c'}) is rather tricky,
because it appears a difference between two infinite quantities,
which has to be `renormalized' in the convenient way.}:
\begin{eqnarray}
\label{for5175}
|S_{(z)}| &\leq& \epsilon(z) \sqrt{V_{(z-1)}W_{(z+1)}} ; \\
\label{for5246}
W_{(z)} &\leq& W_{(z+1)} - S_{(z)}^2 \mathbin{\big/} V_{(z-1)} ; \\
\label{for3534}
V_{(z-1)} &\leq& v - \Big( \sum_{z'<z} S_{(z')} \Big)^{\!2} \mathbin{\Big/} W_{(z)} .
\end{eqnarray}
\end{itemize}
\end{Def}

With that vocabulary, our informal statement can be made precise:
let~$\mbf{A}$ be a compact correct array with entries $V_i^j, W_j^i, S_{ij}$,
and associated quantities $V, W, S$,
and define the arrays $\mbf{A}_k$ by~(\ref{for3010}).
Then when~$k \longto \infty$ one has $(2k+1) \mbf{A}_k \longto \bar{\mbf{A}}$ (in the sense that
each entry of~$(2k+1) \mbf{A}_k$ converges to the corresponding entry of~$\bar{\mbf{A}}$),
where $\bar{\mbf{A}}$ is the Toeplitz array with entries $\bar{V}_i^j,\bar{W}_j^i,\bar{S}_{ij}$
defined by:
\begin{eqnarray}
\bar{V}_{(z)} &=& \sum_{j-i=z} V_i^j ; \\ 
\bar{W}_{(z)} &=& \sum_{j-i=z} W_j^i ; \\
\bar{S}_{(z)} &=& \sum_{j-i=z} S_{ij}.
\end{eqnarray}
%
This array $\bar{\mbf{A}}$ is Toeplitz compact
with $z^-=\min\dot{J}-\max\dot{I}$, resp.\ $z^+=\max\dot{J}-\min\dot{I}$,
and the quantities~(\ref{for3162a})--(\ref{for3162c}) for~$\bar{\mbf{A}}$
are:
\begin{eqnarray}
\bar{v} &=& V ; \\
\bar{w} &=& W ; \\
\bar{s} &=& S.
\end{eqnarray}
Moreover $\bar{\mbf{A}}$ is correct, because all the~$(2k+1) \mbf{A}_k$ are,
and being correct is clearly conserved by array convergence.

The consequence of this statement is the following claim,
which achieves the reduction to a `$\ZZ$-indexed' problem
I alluded to a few lines above:
%
\begin{Clm}
The supremum of~$|S|/\sqrt{VW}$ for correct arrays
is not greater than the supremum of~$|s|/\sqrt{vw}$
for correct Toeplitz arrays.
\end{Clm}

So we have to study (compact) correct Toeplitz arrays.
Consider such an array.
Denote $\theta(z) \coloneqq \arcsin \epsilon(z)$;
then~(\ref{for5175}) can be rewritten:
\[\label{f7312} \exists \hat{\theta}(z) \in [\pm\theta(z)] \qquad
S_{(z)} = \sin \hat{\theta}(z) \cdot \sqrt{V_{(z-1)}W_{(z+1)}} .\]
%
Now, notice that for fixed values of the~$V_{(z)}$, the~$S_{(z)}$ and~$w$,
if we have values $W_{(z)}$ such that (\ref{for5175})--(\ref{for3534}) are satisfied,
we can modify those $W_{(z)}$ so that (\ref{for5246}) becomes an equality for all~$z$,
an operation which keeps (\ref{for5175}) and~(\ref{for3534}) true
since it can only make the~$W_{(z)}$ increase.
So we can suppose that (\ref{for5246}) actually is an equality,
i.e.\ that for all~$z\in\ZZ$,
\[\label{for5357} W_{(z)} = w \prod_{z'\geq z} \cos^2\hat{\theta}(z') .\]
%
Then it remains to integrate~(\ref{for3534}).
For~$z\in\ZZ$, denote
\[ \Gamma(z) \coloneqq \sum_{z'<z} \Big( \sin \hat{\theta}(z') \cdot
\prod_{z'<z''<z} \cos \hat{\theta}(z'') \cdot \sqrt{V_{(z'-1)}} \Big) ,\]
so that (\ref{for3534}) becomes:
\[\label{for6008} V_{(z-1)} \leq v - \Gamma(z)^2 .\]
$\Gamma(\Bcdot)$ satisfies the recursion equation
\[ \Gamma(z+1) = \sin\hat{\theta}(z) \sqrt{V_{(z-1)}} + \cos\hat{\theta}(z) \Gamma(z) ,\]
so by~(\ref{for6008}):
\[\label{for6399} |\Gamma(z+1)| \leq \sin |\hat{\theta}(z)| \sqrt{v-\Gamma(z)^2}
+ \cos\hat{\theta}(z) |\Gamma(z)| .\]
%
From~(\ref{for6399}), we will now prove that for all~$z\in\ZZ$:
\[\label{for6450} |\Gamma(z)| \leq
\sin \Big( \frac{\pi}{2} \wedge \sum_{z'<z} \theta(z') \Big) \sqrt{v} .\]
%
Indeed, (\ref{for6450}) is equivalent to saying that
there exists some $\eta(z) \in [0, \sum_{z'<z} \theta(z')]$
such that $|\Gamma(z)| = \sin\eta(z) \sqrt{v}$,
which we prove by induction.
First, since our Toeplitz array was supposed compact,
$\forall z<z^- \enskip \hat\theta(z) = 0$, so
the formula is true for $z\leq z^-$ with $\eta(z) = 0$.
Next, if the formula is true for~$z$, then (\ref{for6399}) yields
\[|\Gamma(z)| \leq \big( \sin|\hat{\theta}(z)| \cos\eta(z)
+ \cos|\hat{\theta}(z)| \sin \eta(z) \big) \sqrt{v}
= \sin\big( \eta(z) + |\hat{\theta}(z)| \big) \sqrt{v} ,\]
where $\eta(z) + |\hat{\theta}(z)| \leq \sum_{z'<z} \theta(z') + \theta(z)
= \sum_{z'<z+1} \theta(z')$, so the formula is true for~$(z+1)$,
which ends the induction.

To conclude, we write that $s = \sum_z S_{(z)} = \Gamma({z>z^+})\sqrt{w}$.
But by~(\ref{for6450}),
${|\Gamma(z > z^+)|} \ab \leq \sin\bar{\epsilon} \cdot \sqrt{v}$,
so in the end:
\[\label{for4760} |s| \leq \sin\bar{\epsilon} \cdot \sqrt{vw} ,\]
\emph{quod erat demonstrandum.}
\end{proof}


\begin{Cor}[`$\ZZ^n$ against $\ZZ^n$' theorem]\label{corZnZn}
Let~$n\geq 1$; let~$(X_x)_{x\in\ZZ^n}$ and~$(Y_y)_{y\in\ZZ^n}$ be random variables,
and assume there exists a function~$\epsilon \colon \ZZ^n \to [0,1]$
such that for all~$x,y\in\ZZ^n$,
\[ \{ X_x : Y_y \}_{\mcal{M}} \leq \epsilon(y-x) ,\]
$\mcal{M}$ being the natural $\sigma$-metalgebra of the system.
Then $\{\vec{X} : \vec{Y}\} \leq \bar{\epsilon}$,
where $\bar{\epsilon}$ the number in~$[0,1]$ such that
\[\label{for7131} \arcsin (\bar{\epsilon}) =
\Big( \sum_{v\in\ZZ^n} \arcsin \epsilon(v) \Big) \wedge \frac{\pi}{2} .\]
\end{Cor}

\begin{proof}
To alleviate notation, we define the `arcsin-sum'
as the binary operation $\tilde{+} \colon [0,1]^2 \to [0,1]$
defined by:
\[ a \tilde{+} b = \sin \Big( \big( \arcsin a + \arcsin b \big) \wedge \frac{\pi}{2} \Big) .\]
$\tilde{+}$ is associative, commutative and nondecreasing,
so it can be extended into an $\infty$-ary operator $\tilde{\sum}$;
with this notation, (\ref{for7131}) merely writes
$\bar{\epsilon} = \tilde{\sum}_{v\in\ZZ^n} \epsilon(v)$.

Let~$(\mbf{e}_1,\ldots,\mbf{e}_n)$ be a $\ZZ$-basis of~$\ZZ^n$.
For~$1\leq r\leq n$,
we identify $\ZZ^r$ with~$\ZZ\mbf{e}_1 \oplus \ZZ\mbf{e}_2 \oplus \cdots \oplus \ZZ\mbf{e}_r$;
we also denote $\ZZ_r^\perp \coloneqq \ZZ\mbf{e}_{r+1} \oplus \cdots \oplus \ZZ\mbf{e}_n$.
What we will prove is actually the following
%
\begin{Clm}\label{clm6369}
For all~$r\in\{1,\ldots,n\}$, all~$x,y\in\ZZ_r^\perp$,
\[\label{for7222} \big\{ \vec{X}_{x+\ZZ^r} : \vec{Y}_{y+\ZZ^r} \big\}_{\mcal{M}}
\leq \tilde{\sum}_{v\in\ZZ^r} \epsilon(y-x+v) .\]
\end{Clm}%
%
\noindent The statement of the lemma then corresponds to the claim for $r=n$.

We prove Claim~\ref{clm6369} by induction on~$r$.
The case $r=1$ is merely Theorem~\ref{thm0119}%
\footnote{More precisely, it is the subjective version of that theorem,
cf.\ \S~\ref{parSubjectiveResults}.}.
Now let us show how to go from the case $r-1$ to the case $r$ for~$r>1$:

Take $x,y \in \ZZ_r^\perp$.
We notice that
\[ \vec{X}_{x+\ZZ^r} = \big( \vec{X}_{x+i\mbf{e}_r+\ZZ^{r-1}} \big)_{i\in\ZZ} ,\]
which we shorthand into $\vec{X}_{x+\ZZ^r} = (\mbf{X}_i)_{i\in\ZZ}$;
similarly we write, with obvious notation, $\vec{Y}_{y+\ZZ^r} \eqqcolon (\mbf{Y}_j)_{j\in\ZZ}$.
%
By induction hypothesis one has for all~$i,j\in\ZZ$:
\[\label{for7020} \{ \mbf{X}_i : \mbf{Y}_j \}_{\mcal{M}} \leq
\tilde{\sum}_{v\in\ZZ^{r-1}} \epsilon\big(y-x+(j-i)\mbf{e}_r+v\big) .\]
Since the right-hand side of~(\ref{for7020}) only depends on~$(j-i)$,
we can apply Theorem~\ref{thm0119}
to the~$\mbf{X}_i$ and the~$\mbf{Y}_j$, which yields
\[ \big\{ \vec{X}_{x+\ZZ^r} : \vec{Y}_{y+\ZZ^r} \big\}_{\mcal{M}} \leq
\tilde{\sum}_{z\in\ZZ} \Big( \tilde{\sum}_{v\in\ZZ^{r-1}} \epsilon\big(y-x+z\mbf{e}_r+v \big) \Big)
= \tilde{\sum}_{v\in\ZZ^r} \epsilon(y-x+v) ,\]
i.e.~(\ref{for7222}).
\end{proof}




\section{Generalizations of the tensorization results}

\subsection{Minimal Hypotheses}\label{parMinimalHypotheses}

When reading the proofs of the tensorization theorems,
you may have noticed that taking the decorrelation hypotheses
w.r.t.\ the whole $\sigma$-metalgebra of the system was a needlessly strong assumption.
Actually each decorrelation hypothesis can be stated
relatively to only \emph{one} $\sigma$-algebra,
in the following way:
%
\begin{itemize}
\item For Theorem~\ref{thm5750},
one needs only assume that for all~$i\in I$,
$X_i$ and~$Y$ are $\epsilon_i$-decorrelated
when seen from $\sigma\big((X_{i'})_{i'<i}\big)$;
%
\item For Theorems~\ref{thm6413} and~\ref{thm0119},
one needs only assume that
$X_i$ and~$Y_j$ are $\epsilon_{ij}$-decorrelated (or $\epsilon(j-i)$-decorrelated)
when seen from $\sigma\big((X_{i'})_{i'<i},(Y_{j'})_{j'<j}\big)$.
\end{itemize}

In practice it is rare that one can bound above
${\{X_i:Y\}_{\vec{X}_{\{i'<i\}}}}$
or~${\{X_i:Y_j\}_{\mathlarger{(}\vec{X}_{\{i'<i\}},\vec{Y}_{\{j'<j\}}\mathlarger{)}}}$
more sharply than~${\{X:Y_i\}}_{\mcal{M}}$, resp.~${\{X_i:Y_j\}_{\mcal{M}}}$;
yet it is worth remembering that the `genuine' decorrelation hypotheses
are weaker than those we wrote,
especially when one gets interested in optimality issues (cf.\ \S~\ref{parOptimality}).

\begin{Rmk}
In our tensorization proofs we took~$I$ and~$J$ finite; yet those proofs,
and therefore everything in this subsection,
remain valid if we take for~$I$ or~$J$ any (countable) well-ordered set,
in particular if $I$ or~$J$ is~$\NN$.
\end{Rmk}


\subsection{Subjective versions of the theorems}\label{parSubjectiveResults}

In the tensorization theorems I stated,
the decorrelation hypotheses were given with regard to
the natural $\sigma$-metalgebra $\mcal{M}$ of the system,
while the results were given in terms of `objective' (I mean, not subjective) decorrelations.
Yet actually it can be shown that our results are still valid w.r.t.~$\mcal{M}$
—or even w.r.t.\ any sharper $\sigma$-metalgebra $\mcal{N}\supset\mcal{M}$,
provided decorrelation hypotheses are stated w.r.t.~$\mcal{N}$.
As an example, let us state and prove the subjective result corresponding to Theorem~\ref{thm5750}:
%
\begin{Cor}\label{cor4274}
Let~$X$, $(Y_i)_{i\in I}$ and~$(Z_\theta)_{\theta\in\Theta}$ be random variables,
and call~$\mcal{N}$ the $\sigma$-metalgebra they span.
Suppose we have bounds ${\{X:Y_i\}}_{\mcal{N}} \leq \epsilon_i$ for all~$i\in I$;
then:
\[\label{for4193}
\{ X : \vec{Y}_I \}_{\mcal{N}} \leq \sqrt{1-\prod_{i\in I}(1-\epsilon_i^2)} .\]
\end{Cor}

\begin{proof}
Up to making up copies of~$I$ and~$\Theta$,
we can assume that $\{0\}$, $I$ and~$\Theta$ are disjoint,
which allows us to denote $Z_0 \coloneqq X$ and $Z_i \coloneqq Y_i$ for~$i\in I$,
so that $\mcal{N}$ is the $\sigma$-metalgebra spanned by the~$Z_\theta$
for~$\theta\in\bar{\Theta} \coloneqq \{0\} \uplus I \uplus \Theta$.
%
Then (\ref{for4193}) means that for all~$\Xi \subset \bar{\Theta}$,
for (almost-)all~$\vec{z}_{\Xi}$, one must have:
\[ \{ X : \vec{Y}_I \} \leq \sqrt{1-\prod_{i\in I}(1-\epsilon_i^2)} 
\quad \text{under the law $\Pr\big[\Bcdot\big|\vec{Z}_{\Xi} = \vec{z}_{\Xi}\big]$.} \]
%
So, Corollary~\ref{cor4274} will ensue from Theorem~\ref{thm5750}
provided we can prove that, denoting by~$\mcal{M}$ the $\sigma$-metalgebra spanned by~$X$ and the~$Y_i$,
one has for all~$i\in I$:
\[\label{for4395} \{ X : Y_i \}_{\mcal{M}} \leq \epsilon_i
\quad \text{under the law $\Pr\big[\Bcdot\big|\vec{Z}_{\Xi} = \vec{z}_{\Xi}\big]$.} \]

But under a law $P$, saying that ${\{X:Y_i\}}_{\mcal{M}} \leq \epsilon_i$ means that for all~$\Upsilon\subset\{0\}\uplus I$,
for (almost-)all~$\vec{z}\,'_{\!\Upsilon}$,
one has ${\{X:Y_i\}}\leq\epsilon_i$ under the law $P[\Bcdot|\vec{Z}_{\Upsilon} = \vec{z}\,'_{\!\Upsilon}]$.
So, for $P = \Pr[\Bcdot\big|\vec{Z}_{\Xi} = \vec{z}_{\Xi}]$,
(\ref{for4395}) means that, for all~$\vec{z}\,'_{\!\Upsilon}$:
%
\[\label{for4586} \{ X : Y_i \} \leq \epsilon_i
\quad \text{under the law $\Pr\big[\Bcdot\big|\vec{Z}_{\Xi} = \vec{z}_{\Xi}\enskip\text{and}\enskip \vec{Z}_{\Upsilon} = \vec{z}\,'_{\!\Upsilon}\big]$.} \]
In Formula~(\ref{for4586}) we can assume that $z_\theta$ and~$z'_\theta$ coincide
for all~$\theta \in \Xi \cap \Upsilon$,
since otherwise the event ``$\vec{Z}_{\Xi}=\vec{z}_{\Xi}\enskip\text{and}\enskip \vec{Z}_{\Upsilon} = \vec{z}\,'_{\!\Upsilon}$''
would be empty and there would be nothing to say.
Then ``$\vec{Z}_{\Xi}=\vec{z}_{\Xi}\enskip\text{and}\enskip \vec{Z}_{\Upsilon} = \vec{z}\,'_{\!\Upsilon}$''
is of the form ``$\vec{Z}_{\Xi\cup\Upsilon} = \vec{z}_{\Xi\cup\Upsilon}$'',
where $\Xi\cup\Upsilon \subset \bar{\Theta}$,
so that (\ref{for4586}) follows directly from the hypothesis
$\{X:Y_I\}_{\mcal{N}} \leq \epsilon_i$.
\end{proof}




\section{Optimality}\label{parOptimality}

\subsection{Exact Optimality}

With the minimal hypotheses stated in \S~\ref{parMinimalHypotheses},
Theorems~\ref{thm5750} and~\ref{thm0119} are optimal:
%
\begin{Thm}\label{t6656}
The bound~(\ref{for5719}) in Theorem~\ref{thm5750} is optimal, in the following sense:
for any integer $N$, for all~$(\epsilon_i)_{1\leq i\leq N}$ in~$[0,1]^N$,
one can find random variables $X_1, \ldots, X_N, Y$ such that
for all~$i \in \{1,\ldots,N\}$,
\[\label{f6808} \{X_i:Y\}_{\vec{X}_{\{i'<i\}}} = \epsilon_i \]
and
\[ \{ \vec{X} : Y \} = \sqrt{1-\prod_{i}(1-\epsilon_i^2)} .\]
\end{Thm}

\begin{Thm}\label{thm6657-0}
The bound~(\ref{f6444}) in Theorem~\ref{thm0119} is optimal, in the following sense:
for any integer $N$, for all~$(\epsilon(z))_{-N\leq z\leq N} \in [0,1]^{\{-N,\ldots,N\}}$,
one can find random variables~$(X_i)_{i\in\ZZ}$ and~$(Y_j)_{j\in\ZZ}$ such that for all~$i,j\in\ZZ$,
\[\label{f0794} \{ X_i : Y_j \}_{\mathlarger{(}\vec{X}_{\{i'<i\}},\vec{Y}_{\{j'<j\}}\mathlarger{)}}
= \left\{ \begin{array}{lcl} \epsilon(j-i) & \quad & \text{if $|j-i| \leq N$;} \\
0 & \quad & \text{if $|j-i| > N$} \end{array} \right.\]
and ${\{ \vec{X} : \vec{Y} \}} = \bar{\epsilon}$, with~$\bar{\epsilon}$ defined by:
\[\label{f6595} \arcsin \bar{\epsilon} = \sum_{z=-N}^N \arcsin \epsilon(z)
\wedge \frac{\pi}{2} .\]
\end{Thm}

Actually, as proving Theorem~\ref{thm6657-0} for \emph{all} the~$(\epsilon(z))_{-N\leq z\leq N}$
involves some heavy technicalities~\cite{perso},
I will only prove the slightly weaker following
%
\begin{Thm}\label{t6657}
For any integer $N$, the exists a neighbourhood $U$ of~$\vec{0}$ in~$[0,1]^{\{-N,\ldots,N\}}$
such that, for all~$(\epsilon(z))_{-N\leq z\leq N} \in U$, one can find random variables
$(X_i)_{i\in\ZZ}$ and~$(Y_j)_{j\in\ZZ}$ satisfying~(\ref{f0794}) and~(\ref{f6595})%
\footnote{Notice that in the neighbourhood of~$\vec{0}$,
one can drop the ``$\wedge\frac{\pi}{2}$'' in the right-hand side of~(\ref{f6595}).}.
\end{Thm}

\begin{Rmk}
On the other hand, Theorem~\ref{thm6413}
is obviously not optimal since, as we pointed out, its bound is strictly weaker
than that of Theorem~\ref{thm0119}.
\end{Rmk}

The proof of Theorem~\ref{t6656} relies on the following important result:
%
\begin{Lem}\label{l6791}
Let~$(X_1,\ldots,X_N,Y)$ be an $(N+1)$-dimensional Gaussian vector.
For all~$1\leq i\leq N$, define
\[ e_i \coloneqq \{ X_i : Y \}_{\vec{X}_{\{i'<i\}}} ,\]
then one has exactly:
\[\label{f1922} \{ \vec{X} : Y \} = \sqrt{1-\prod_{i}(1-e_i^2)} .\]
\end{Lem}

\begin{Rmk}\label{rmk>3lines}
Maximal correlation, as I told in \S~\ref{parBasicHilbertianDecorrelations},
is fundamentally a Hilbertian concept. When one deals with Gaussian vectors,
the Hilbert spaces involved actually have finite dimensions,
so that Lemma~\ref{l6791} about decorrelations can also be seen as a result about Euclidian spaces.
In Appendix~\ref{par3lines}, I will present an unexpected corollary of this lemma,
stating a geometric property of the $3$-dimensional Euclidian space.
\end{Rmk}

\begingroup\def\proofname{Proof of Lemma~\ref{l6791}}\begin{proof}
To alleviate notation, we denote $\mcal{F}_{i-1} \coloneqq \sigma\big(\vec{X}_{\{i'<i\}}\big)$.
Since $(\vec{X},Y)$ is Gaussian, the law of~$(X_i,Y)$
under~$\Pr[\Bcdot|\ab x_1,\ldots,x_{i-1}]$ is Gaussian and only depends on
$(x_1,\ldots,x_{i-1})$ through an additive constant;
%
consequently, we can speak of ``\emph{the} Hilbertian correlation between~$X_i$ and~$Y$ conditionally to~$\mcal{F}_{i-1}$'', which is $e_i$,
and also of ``\emph{the} conditional variance of~$X_i$ w.r.t.~$\mcal{F}_{i-1}$'',
resp.\ ``\emph{the} conditional variance of~$Y$'', resp.\ ``\emph{the} conditional covariance of~$(X_i,Y)$'',
which we denote resp.~$\Var(X_i|\mcal{F}_{i-1})$, $\Var(Y|\mcal{F}_{i-1})$,
$\Cov(X_i,Y|\mcal{F}_{i-1})$. By Theorem~\ref{pro1857}, one has:
\[\label{f2664}
\Cov(X_i,Y|\mcal{F}_{i-1}) = \pm e_i \ecty(X_i|\mcal{F}_{i-1}) \ecty(Y|\mcal{F}_{i-1}) .\]

Now take $g(Y) = Y$ and $f(X) = \sum_{i=1}^{N} \beta_i X_i$,
for some $\beta_i \in \RR$ to be chosen later.
Then $g^{i-1}$ is equal to~$Y - \EE[Y|\mcal{F}_{i-1}]$ and $f_i$ is proportional
to~$X_i - \EE[X_i|\mcal{F}_{i-1}]$, thus, by~(\ref{f2664}) and our model's being Gaussian,
all the inequalities until~(\ref{f9041c}) in the proof of Theorem~\ref{thm5750}
actually are equalities for $\epsilon_i = e_i$.
If moreover $\Cov(f_i,g^{i-1}|\mcal{F}_{i-1}) \geq 0$ for all~$i$,
then we can drop the absolute values in their left-hand sides,
and thus (\ref{f8485}) will also be an equality.
Then, to get an equality in~(\ref{f9280}), it just remains to ensure
that the final Cauchy--Schwarz equality is an equality,
i.e.\ to ensure that one has, for all~$i$:
\[\label{f2464} \Var(f_i) \propto e_i^2 \prod_{i'=1}^{i-1} (1-e_{i'}^2) .\]
%
If all of that is satisfied, then one will have exactly
$\EE[fg] = \sqrt{1-\prod_{i}(1-e_i^2)} \ab \hspace{.056em} \ecty(f)\ecty(g)$,
so that $\{\vec{X}:Y\} \geq \sqrt{1-\prod_{i}(1-e_i^2)}$.
The converse inequality being obviously true
by (the minimal version of) Theorem~\ref{thm5750},
the result will follow.

So, we have to check that the choice of the~$\beta_i$ can be performed so that
(\ref{f2464}) is satisfied, with $\Cov(f_i,g^{i-1}|\mcal{F}_{i-1})$ of the good sign.
To do this, we will choose successively
relevant values for~$\beta_N,\ab \beta_{N-1},\ab \ldots,\ab \beta_1$.

We observe that, if $\beta_N,\ldots,\beta_{i+1}$ have already been fixed, then
$\beta_i \mapsto \Cov(f_i,g^{i-1}|\mcal{F}_{i-1})$
is an affine function with slope
\[ \pm e_i \ecty(Y|\mcal{F}_{i-1}) \frac{\ecty(X_i|\mcal{F}_{i-1})}{\ecty(X_i)} .\]
Moreover, $\Var(f_i) = \Var(f_i|\mcal{F}_{i-1})$ as $f_i$ is centered w.r.t.~$\mcal{F}_{i-1}$;
so, since $f_i \propto X_i - \EE[X_i|\mcal{F}_{i-1}]$, (\ref{f2664}) implies:
\[ \Var(f_i) = \frac{\Cov(f_i,g^{i-1}|\mcal{F}_{i-1})^2}{e_i^2 \Var(Y|\mcal{F}_{i-1})} .\]
%
So, provided all the three quantities $e_i$, $\ecty(Y|\mcal{F}_{i-1})$ and~$\ecty(X_i|\mcal{F}_{i-1})$ are nonzero,
there exists a (unique) $\beta_i$ satisfying~(\ref{f2464}).

Now if $\ecty(Y|\mcal{F}_{i-1})$ is zero, this means that $Y$ is $\mcal{F}_{i-1}$-measurable;
then one of the~$e_{i'}$ has to be~$1$ and thus the result is trivial.
Next if $\ecty(X_i|\mcal{F}_{i-1})$ is zero,
this means that $X_i$ is $\mcal{F}_{i-1}$-measurable;
then $e_i = 0$ and $f_i \equiv 0$, so that (\ref{f2464}) is automatically satisfied.
Finally if $e_i = 0$ and $\Var(X_i|\mcal{F}_{i-1})>0$, then there exists a (unique) $\beta_i$
such that $f_i \equiv 0$, for which (\ref{f2464}) is satisfied.
So all those particular cases actually work fine too.
\end{proof}\endgroup

\begingroup\def\proofname{Proof of Theorem~\ref{t6656}}\begin{proof}
For technical reasons, we begin with noticing that the theorem
is immediate if some $e_i$ is equal to~$1$,
so that we can assume that all the~$e_i$ are $<1$.
Thanks to Lemma~\ref{l6791}, it suffices to prove that
for any sequence of~$\epsilon_i \in [0,1)$
it is possible to build a Gaussian vector~$(X,\vec{Y})$
for which $e_i = \epsilon_i \ \forall i$.
To do this, let~$\xi, \zeta_1, \ldots, \zeta_N$ be i.i.d.\ $\mcal{N}(1)$ variables,
and take $Y = \xi$ and $X_i = \sqrt{1-\alpha_i}\zeta_i + \sqrt{\alpha_i}\xi$
for some parameters $\alpha_i \in [0,1)$.
We want to choose the~$\alpha_i$ such that $\vec{e}(\vec{\alpha}) = \vec{\epsilon}$;
this is always possible, by the following method:
\begin{itemize}
\item First we compute $\alpha_1$: By Theorem~\ref{pro1857},
one can write down the equation linking~$\alpha_1$ and~$e_1$.
It is clear without knowing the precise form of that equation
(actually, $e_1 = \sqrt{\alpha_1}$)
that $e_1$ is a continuous increasing function of~$\alpha_1$
with $e_1=0$ for $\alpha_1=0$ and $e_1=1$ for $\alpha_1=1$.
Therefore there is a unique $\alpha_1$ such that $e_1=\epsilon_1$.
\item Then we compute $\alpha_2$: As we already know the value of~$\alpha_1$,
we can treat it as a constant and look for the equation linking~%
$\alpha_2$ and~$e_2$, which we compute by Theorem~\ref{pro1857} again.
That equation, though more complicated than in the previous case
(actually, $e_2 = \sqrt{\alpha_2}\sqrt{1-\alpha_1}/\sqrt{1-\alpha_1\alpha_2}$),
exhibits the same behaviour: $e_2$ is a continuous increasing function of~$\alpha_2$
with $e_2(\alpha_2=0)=0$ and $e_2(\alpha_2=1)=1$.
Therefore there is a unique $\alpha_2$ such that $e_2=\epsilon_2$.
\item We carry on this process until having determined all the~$\alpha_i$.
\end{itemize}
\end{proof}\endgroup


\begingroup\def\proofname{Proof of Theorem~\ref{t6657}}\begin{proof}
Again, the principle of the proof will consist in showing
how the optimal bound can be attained for relevant Gaussian vectors
and linear functions of them.

We consider independent $\mcal{N}(1)$ variables~$(\xi_j)_{j\in\ZZ}$
and~$(\omega_{ij})_{(i,j)\in\ZZ\times\ZZ}$.
For all~$i$ we set:
\[ X_i = \sum_{z=-N}^{N} \omega_{i(i+z)} ,\]
resp. for all~$j$:
\[ Y_j = \xi_j + \sum_{z=-N}^{N} \alpha_z \omega_{(j-z)j} \]
for some real parameters $(\alpha_z)_{-N\leq z\leq N}$ to be fixed later.
This model is obviously invariant by translation of the indexes.
For~$z\in\ZZ$, define
\[ \dot{e}_z \coloneqq \frac{\Cov\big(X_i,Y_{i+z}\big|\mcal{F}_{i-1}\vee\mcal{G}_{i+z-1}\big)}
{\ecty\big(X_i\big|\mcal{F}_{i-1}\vee\mcal{G}_{i+z-1}\big)
\, \ecty\big(Y_{i+z}\big|\mcal{F}_{i-1}\vee\mcal{G}_{i+z-1}\big)} ,\]
where the choice of~$i$ does not matter.
Since our model is Gaussian, by Theorem~\ref{pro1857},
\[ \{ X_i : Y_{i+z} \}_{\mathlarger{(}\vec{X}_{\{i'<i\}},\vec{Y}_{\{j'<i+z\}}\mathlarger{)}}
= |\dot{e}_z| .\]

By the properties of Gaussian vectors,
it is possible to write down explicitly the equations linking
the~$\dot{e}_z$ to the~$\alpha_z$.
Though these equations may be quite horrendous, some of their properties can be easily established:
\begin{Clm}\label{c4330}\strut
\begin{ienumerate}
\item\label{i4332} For $|z|>N$, $\dot{e}_z = 0$ (for any choice of the~$\alpha_z$);
\item The map $(\alpha_{-N},\ldots,\alpha_N) \mapsto (\dot{e}_{-N},\ldots,\dot{e}_N)$
is of class $\mcal{C}^1$ on the neighbourhood of~$(0,\ldots,0)$, with:
\[ \left(\frac{\partial \dot{e}_z}{\partial \alpha_y}\right)(\vec0) =
\frac{\1{y=z}}{\sqrt{2N+1}} .\]
\end{ienumerate}
\end{Clm}
%
\noindent By the inverse function theorem, one can therefore find neighbourhoods~%
$V$ and~$U$ of~$\vec0$ in~$\RR^{\{-N,\ldots,N\}}$ such that
the map $\vec{\alpha} \mapsto \vec{\dot{e}}$ is a $\mcal{C}^1$-diffeomorphism from $V$ onto $U$.
In particular, for~$\vec\epsilon$ in such an $U$ we can always fix the~$\alpha_z$
of our model such that $\forall z\enskip\dot{e}_z = \1{|z|\leq N}\epsilon(z)$,
so that (\ref{f0794}) is satisfied.

Now we have to choose~$f$ and~$g$.
Morally\footnote{I say ``morally'' because nothing ensures
that the supremum~(\ref{for3878}) would actually be a maximum here.}
we have to take the functions~$f$ and~$g$ having maximal Pearson correlation.
Since the model is Gaussian, these functions will be linear,
and since the model is invariant by translation,
they will likely be invariant by translation too.
So we would like to take, formally, $f(\vec{X}) = \sum_{i\in\ZZ} X_i$
and $g(\vec{Y}) = \sum_{j\in\ZZ} Y_j$.
%
As such functions are not properly defined,
we will rather consider
$f[k](\vec{X}) = \sum_{i=-k}^k X_i$,
resp.\ $g[k](\vec{Y}) = \sum_{j=-k}^k Y_j$,
and then we will let~$k$ tend to infinity.

For these~$f[k]$ and~$g[k]$, define the~$V[k]^j_i$, the~$W[k]_j^i$ and the~$S[k]_{ij}$
as in the proof of Theorem~\ref{thm0119}, which are gathered into the array $\mbf{A}[k]$.
%
The following properties of the~$\mbf{A}[k]$
follow easily from the structure of our model:
%
\begin{Clm}\label{c6626}\strut
\begin{ienumerate}
\item All the~$V[k]^j_i, W[k]_j^i, S[k]_{ij}$ are bounded uniformly in~$i,j,k$.
\item\begin{itemize}\item $V[k]^j_i$ is zero as soon as $i\notin\{-k-2N,\ldots,k\}$;
\item $W[k]_j^i$ is zero as soon as $j\notin\{-k,\ldots,k\}$.\end{itemize}
\item $S[k]_{ij}$ is zero as soon as $|j-i| > N$.
\item\label{itmHOT1}\begin{itemize}\item For $-k\leq i\leq k-2N$,
$V[k]^j_i$ only depends on~$(j-i)$, even when $k$ varies. We denote its value by~$V_{(j-i)}$.
\item For $-k\leq j\leq k$,
$W[k]_j^i$ only depends on~$(j-i)$, even when $k$ varies. We denote its value by~$W_{(j-i)}$.
\item For $-k\leq i\leq k-2N$ and $-k\leq j\leq k$,
$S[k]_{ij}$ only depends on~$(j-i)$, even when $k$ varies. We denote its value by~$S_{(j-i)}$.\end{itemize}
\item\label{itmHOT2}\begin{itemize}\item $V_{(z)}$ has some constant value $v$ for $z < -N$;
\item $W_{(z)}$ has some constant value $w$ for $z > N$.
\end{itemize}
\end{ienumerate}
\end{Clm}
%
\noindent By Claim~\ref{c6626}, $\mbf{A}[k]$ converges pointwise to
some compact Toeplitz array $\mbf{A}$, whose entries are the~$V_{(z)}, W_{(z)}, S_{(z)}$
introduced at Item~(\ref{itmHOT1}) of the claim,
whose values~$v$ and~$w$ are those introduced at Item~(\ref{itmHOT2}),
and whose value $s$ is $\sum_{z=-N}^{N} S_{(z)}$.
All the arrays $\mbf{A}[k]$ are obviously correct since they correspond to true functions,
so by passing to the limit $\mbf{A}$ is correct too.

Since our model is Gaussian,
all the inequalities~(\ref{for9472a}), (\ref{for9472b}) and~(\ref{for9472c})
are actually equalities for the arrays $\mbf{A}[k]$;
moreover, since the~$\dot\epsilon_z$ are nonnegative,
the~$S[k]_{ij}$ are nonnegative.
By letting $k$ tend to infinity, it follows that
all the inequalities (\ref{for5175})--(\ref{for3534}) are actually equalities for the array $\mbf{A}$,
with the~$S_{(z)}$ nonnegative.
Consequently in~(\ref{f7312}) one has $\hat\theta(z) = \theta(z)$,
and all the further inequalities are actually equalities, so that in the end
(\ref{for4760}) becomes:
\[\label{for4760=} \frac{s}{\sqrt{vw}} = \bar\epsilon .\]

Now, defining~$V[k]$, $W[k]$ and~$S[k]$
by resp.~(\ref{for0923a}), (\ref{for0923b}) and~(\ref{for5163}) for the arrays $\mbf{A}[k]$,
Claim~\ref{c6626} shows that, when~$k \longto \infty$,
$V[k] \sim 2kv$, resp.\ $W[k] \sim 2kw$, resp.\ $S[k] \sim 2ks$,
so~(\ref{for4760=}) implies that $S[k]/\sqrt{V[k]W[k]} \longto \bar\epsilon$.
But recall that $V[k]$, $W[k]$ and~$S[k]$ are the respective variances and covariance
of the functions~$f[k] \in \ldb(\vec{X})$ and~$g[k] \in \ldb(\vec{Y})$,
so by the very definition~(\ref{for3878}) of Hilbertian correlations,
\[ \{\vec{X} : \vec{Y}\} \geq \frac{S[k]}{\sqrt{V[k]W[k]}} .\]
Making $k \longto \infty$, it follows that $\{\vec{X}:\vec{Y}\} \geq \bar\epsilon$;
the converse inequality being obviously true
by (the minimal version of) Theorem~\ref{thm0119}, this proves Theorem~\ref{t6657}.
\end{proof}\endgroup

\begin{Xpl}\label{xplOptZZ}
In this example we will carry out explicit computations for a Gaussian model
close to the model presented in the proof above.
%
We take independent $\mcal{N}(1)$ variables $\ldots,\zeta_{-1},\zeta_0,\zeta_1,\ldots$,
$\ldots,\xi_{-1/2},\xi_{1/2},\xi_{3/2}\ldots$,
$\ldots,\omega_{-1/4},\omega_{1/4},\omega_{3/4},\ldots$,
and we set
%
\begin{eqnarray}
X_i &=& \zeta_i + \sqrt{\alpha} (\omega_{i-1/4} + \omega_{i+1/4}) , \\
\llap{resp.\quad} Y_j &=& \xi_j + \sqrt{\alpha} (\omega_{j-1/4} + \omega_{j+1/4})
\end{eqnarray}
for all integer $i$, resp.\ all half-integer $j$,
where $\alpha$ is some arbitrary nonnegative parameter.
We are going to show that for this system (\ref{f6444}) is actually an equality,
in accordance with the proof of Theorem~\ref{t6657}.

For half-integer $z$ denote
\[ e_z \coloneqq \{ X_i : Y_{i+z} \}%
_{\mathlarger{(}\vec{X}_{\{i'<i\}},\vec{Y}_{\{j'<i+z\}}\mathlarger{)}} ,\]
where the choice of~$i$ does not matter by translation invariance.
Clearly $e_{-z} = e_z$ for all~$z$ and $e_z = 0$ for $|z| > 1/2$,
so to know all the~$e_z$ the only nontrivial computation is computing $e_{1/2}$.
Let us perform it.

Since everything is Gaussian, by Theorem~\ref{pro1857}, $e_{1/2}$ is the value,
under the law $\Pr[\Bcdot|\ab \vec{X}_{\{i<0\}},\vec{Y}_{\{j<1/2\}} \equiv0]$, of
\[\label{forVespr} |\EE[X_0Y_{1/2}]| \mathbin{\big/} \ecty(X_0)\ecty(Y_{1/2}) .\]
%
Under the law $\Pr[\Bcdot|\vec{X}_{\{i<0\}},\vec{Y}_{\{j<1/2\}} \equiv0]$,
it is clear that $\zeta_0,\ab \omega_{1/4},\ab \xi_{1/2},\ab \omega_{3/4},\ldots$
have exactly the same (joint) law as under~$\Pr$,
and that $\omega_{-1/4}$ is still independent of these (joint) variables,
though its variance shall have diminished.
So we need only compute
\[ v \coloneqq \Var\big(\omega_{-1/4}\big|\vec{X}_{\{i<0\}},\vec{Y}_{\{j<1/2\}} \equiv0\big) .\]

Denote $\vec{L}_{\mathrm r} \coloneqq (\ldots,X_{-2},Y_{-3/2},X_{-1},Y_{-1/2})$,
resp.\ $\vec{L}_{\mathrm l} \coloneqq (\ldots,X_{-2},Y_{-3/2},X_{-1})$.
%
We write that (formally)
\[ \dx\Pr\big[ \vec{L}_{\mathrm r} \equiv0 \ \text{and} \ \omega_{-1/4} = x \big]
\propto e^{-x^2/2v} \dx{x}, \]
and also
$\dx\Pr[ \vec{L}_{\mathrm l} \equiv0 \ \text{and} \ \omega_{-3/4} = y]
\propto e^{-y^2/2v} \dx{y}$
by translation invariance.
But under~$\Pr[\Bcdot| \ab {\vec{L}_{\mathrm l} \equiv0} \enskip \text{and} \enskip {\omega_{-3/4} = y}]$,
the law of~$(\xi_{-1/2},\omega_{1/4})$ is the same as under~$\Pr$, so one has:
%
\begin{multline}
e^{-x^2/2v} \propto
\dx\Pr\big[ \vec{L}_{\mathrm r} \equiv0 \ \text{and} \ \omega_{-1\!/\!4} = x \big] \\
= \int_y \dx{y} \,
\dx\Pr\big[ \vec{L}_{\mathrm l} \equiv0 \ \text{and} \ \omega_{-3\!/\!4} = y \big]
\dx\Pr\big[ Y_{-1/2} = 0 \ \text{and} \ \omega_{-1\!/\!4} = x \big|
\vec{L}_{\mathrm l} \equiv0 \ \text{and} \ \omega_{-3\!/\!4} = y \big] \\
\propto \int_y \dx\Pr\big[ Y_{-1/2} = 0 \ \text{and} \ \omega_{-1\!/\!4} = x \big|
\vec{L}_{\mathrm l} \equiv0 \ \text{and} \ \omega_{-3\!/\!4} = y \big] e^{-y^2/2v} \,\dx{y} \\
= \int_y
\dx\Pr\big[ \xi_{-1\!/\!2} = -\sqrt{\alpha}(x+y) \ \text{and} \ \omega_{-1\!/\!4} = x \big|
\vec{L}_{\mathrm l} \equiv0 \ \text{and} \ \omega_{-3\!/\!4} = y \big]
e^{-y^2/2v} \,\dx{y} \\
\propto \int_y e^{-\alpha(x+y)^2/2} e^{-x^2/2} e^{-y^2/2v} \,\dx{y}
\propto \exp \Big\{ \Big( 1+\alpha-\frac{\alpha^2}{\alpha+1/v} \Big) \frac{x^2}{2} \Big\},
\end{multline}
so that $v$ must satisfy:
\[ 1+\alpha-\frac{\alpha^2}{\alpha+1/v} = \frac{1}{v} ,\]
whose only nonnegative solution is
\[ v = \frac{\sqrt{1+4\alpha}-1}{2\alpha} .\]

So one has $\ecty(X_0|\vec{L}_{\mathrm l} \equiv0) = \sqrt{1+\alpha v+\alpha}
= \big(\sqrt{1+4\alpha}+1\big)/2$,
$\ecty(Y_{1/2}|\vec{L}_{\mathrm l} \equiv0) = \sqrt{1+2\alpha}$
and $\EE[X_0Y_{1/2}|\vec{L}_{\mathrm l} \equiv0] = \alpha$,
so that in the end (\ref{forVespr}) yields:
\[\label{f3260} e_{1/2} = \frac{\sqrt{1+4\alpha}-1}{2\sqrt{1+2\alpha}} .\]
%
With this value, Theorem~\ref{thm0119} states that one has necessarily
\[\label{forMncz} \{ \vec{X} : \vec{Y} \} \leq \sin (2\arcsin e_{1/2}) \footnotemark
= 2 e_{1/2} \sqrt{1-e_{1/2}^2} = \frac{2\alpha}{1+2\alpha} .\]
\footnotetext{As here one always has $e_{1/2} \leq 1/\sqrt{2}$,
we can drop the ``$\wedge\frac{\pi}{2}$'' of Formula~(\ref{f6444}).}

We show that (\ref{forMncz}) is actually an equality:
take indeed~$f[k](\vec{X}) \coloneqq \sum_{i=1}^{k} X_k$,
resp.\ $g[k](\vec{Y}) \coloneqq \sum_{j=1/2}^{k-1/2} Y_k$,
then $\Var(f[k]) = \Var(g[k]) = k(1+2\alpha)$
and $\EE[fg] = (2k-1)\alpha$,
so that
\[\label{f4151} \{ \vec{X} : \vec{Y} \} \geq \frac{(2k-1)\alpha}{k(1+2\alpha)} 
\stackrel{k\longto\infty}{\longto} \frac{2\alpha}{1+2\alpha} .\]
\end{Xpl}

\begin{Rmk}\label{rmkPhTr}
One can formally set $\alpha = +\infty$ in the previous example,
which actually means that one takes $X_i = \omega_{i-1/4} + \omega_{i+1/4}$,
resp.\ $Y_j = \omega_{j-1/4} + \omega_{j+1/4}$.
In this case, both Formulas~(\ref{f3260}) and~(\ref{f4151}) `pass to the limit',
yielding $e_{1/2} = 1/\sqrt{2}$ and $\{\vec{X}:\vec{Y}\}=1$.
This shows that it is possible indeed that the~$e_z$ have `mild' values
and that yet $\vec{X}$ and~$\vec{Y}$ are fully correlated.
%
In other words, the ``$\wedge\frac{\pi}{2}$'' in~(\ref{f6444})
is not an `artifact' of the proof of Theorem~\ref{thm0119}%
%
\footnote{\label{ftn6129}%
On the other hand, it is possible that the ``$\wedge1$'' in~(\ref{for3009})
was such an artifact, since Theorem~\ref{thm6413} is not optimal.},
%
but the expression of a real `phase transition' phenomenon%
%
\footnote{There exist indeed situations going
`beyond the phase transition', i.e.\ for which $\sum_{z\in\ZZ} \arcsin(e_z) \ab > \pi/2$,
though this is not the case for Example~\ref{xplOptZZ}.}.
%
Such a phase transition did not occur for the simple tensorization formula~(\ref{for5719}),
which shows that double tensorization in intrinsically more complicated
than simple tensorization.
\end{Rmk}


\subsection{Asymptotic optimality}\label{parAsymptoticOptimality}

In the previous subsection we saw that
(the minimal versions of) Theorems~\ref{thm5750}
and~\ref{thm0119} were optimal, while Theorem~\ref{thm6413} was not.
However it turns out that that result is nevertheless `asymptotically optimal',
in the sense that the bound it gives is equivalent to the optimal bound
when the correlations between the variables become weak.
Here is a precise statement:
%
\begin{Thm}
Let $I = \{1,\ldots,N\}$ and $J=\{1,\ldots,M\}$ be finite sets, and define
the function~$\mathit{Opt} \colon [0,1]^{I\times J} \to [0,1]$ by
\[\label{f5516}
\mathit{Opt} \big( \vec\epsilon_{I\times J} \big)
\coloneqq \sup \Big\{ \big\{ \vec{X}_I : \vec{Y}_J \big\} \ ;\ \big(\forall (i,j) \in I\times J\big) \,
\big(\{ X_i:Y_j \}_{\mathlarger{(}\vec{X}_{\{i'<i\}},\vec{Y}_{\{j'<j\}}\mathlarger{)}}
\leq \epsilon_{ij}\big) \Big\} ;\]
then, when~$\vec\epsilon_{I\times J} \longto \vec{0}$,
one has:
\[\label{f7160}
\mathit{Opt} \big(\vec\epsilon\big) \sim \VERT \boldepsilon \VERT .\]
\end{Thm}

\begin{Rmk}
In the same way, the simple bound~(\ref{for5528}) of Proposition~\ref{pro5527}
is asymptotically equivalent to the optimal bound~(\ref{for5719}) of Theorem~\ref{thm5750}.
\end{Rmk}

\begin{proof}
Take $(M+NM)$ i.i.d.\ $\mcal{N}(1)$ variables
$\xi_1,\ldots,\xi_M,\omega_{11},\ldots,\omega_{NM}$.
For~$(\!(\alpha_{ij})\!)_{i,j} \in \RR^{N\times M}$,
set
\[ \left\{ \begin{array}{rcl} X_i &=& \sum_j \omega_{ij} ; \\
Y_j &=& \xi_j + \sum_i \alpha_{ij} \omega_{ij} .\end{array} \right.\]

Denote
\[ e_{ij} \coloneqq
\{ X_i:Y_j \}_{\mathlarger{(}\vec{X}_{\{i'<i\}},\vec{Y}_{\{j'<j\}}\mathlarger{)}} ,\]
and define $\dot{e}_{ij}$ as the Pearson correlation coefficient
of~$X_i$ and~$Y_j$ under the law
$\Pr[\Bcdot|\ab \vec{X}_{\{i'<i\}}, \ab {\vec{Y}_{\{j'<j\}}\equiv 0}]$.
Then, as in the proof of Theorem~\ref{t6657}, one has $e_{ij} = |\dot{e}_{ij}|$, and the function
$\vec\alpha \mapsto \vec{\dot{e}}$ is $\mcal{C}^1$ around~$\vec{0}$, with
\[\label{f8236} \vec{\dot{e}} = \frac{1}{\sqrt{M}}\,\vec\alpha + O(\|\vec\alpha\|^2)
\qquad \text{when $\vec\alpha \longto \vec0$.} \]
%
By the inverse function theorem,
$\vec\alpha \mapsto \vec{\dot{e}}$ is therefore a diffeomorphism
from some neighbourhood~$V$ of~$\vec{0}$ onto some neighbourhood $U$ of~$\vec{0}$,
whose inverse function is such that
\[ \vec{\alpha} = \sqrt{M}\,\vec{\dot{e}} + O(\|\vec{\dot{e}}\|^2) \qquad
\text{when $\vec{\dot{e}} \longto \vec0$.} \]

Now let~$\vec\epsilon \in (\RR_+)^{N\times M} \cap U$.
Take $\vec\alpha\in V$ such that $\vec{\dot{e}}(\vec\alpha) = \vec\epsilon$,
so that the condition of~(\ref{f5516}) is satisfied.
For~$\phi\in\RR^N,\ab {\psi\in\RR^M}$ with $\|\phi\|,\|\psi\| = 1$, set
\[\left\{\begin{array}{rcl}
f(\vec{X}) &\coloneqq& \sum_i \phi_i X_i ; \\
g(\vec{Y}) &\coloneqq& \sum_j \psi_j Y_j.
\end{array}\right.\]
%
One has
\[ \Var(f) = M ,\]
\[ \Var(g) = 1 + O(\|\vec\alpha\|^2) = 1 + O(\|\vec\epsilon\|^2) \]
and
\[ \EE[fg] = \sum_{i,j} \alpha_{ij}\phi_i\psi_j
= \langle \phi, \boldsymbol\epsilon \psi \rangle + O(\|\vec\epsilon\|^2) ,\]
where the constants implicit in the ``$O(\|\vec\epsilon\|^2)$'' are uniform
in~$(\phi,\psi)$.
%
So one has
\[ \mathit{Opt}(\vec{\epsilon}) \geq \{ \vec{X} : \vec{Y} \}
\geq \frac{|\EE[fg]|}{\ecty(f)\ecty(g)}
= |\langle \phi, \boldsymbol\epsilon \psi \rangle| + O(\|\vec\epsilon\|^2) ,\]
whence after taking supremum over $(\phi,\psi)$:
\[ \mathit{Opt}(\vec{\epsilon}) \geq \VERT \boldsymbol\epsilon \VERT + O(\|\vec\epsilon\|^2)
\stackrel{\vec\epsilon\longto\vec0}{\sim} \VERT \boldsymbol\epsilon \VERT .\]
Since on the other hand $\mathit{Opt}(\vec{\epsilon}) \leq \VERT\boldsymbol\epsilon\VERT$
by Theorem~\ref{thm6413}, the proposition follows.
\end{proof}

\begin{Rmk}
If we state decorrelation hypotheses
w.r.t.\ the \emph{whole} $\sigma$-metalgebra of the system (denoted by~$*$),
no quantity analogous to~$\dot{e}_{ij}$ shall exist any more;
then one can only write, denoting $e'_{ij} \coloneqq \{X_i:Y_j\}_{*}$:
\[ e'_{ij} (\vec\alpha) = \frac{|\alpha_{ij}|}{\sqrt{M}} + O(\|\vec\alpha\|^2) .\]
So, to see how the correlations depend on the parameters,
we have to study the map $\vec\alpha\mapsto\vec{e}\,'$,
which is approximated by a homothety \emph{only on the cone $\RR^{N\times M}_+$}
—and which moreover is no better than continuous here.
So we shall replace the inverse function theorem by an alternative technique,
which will yield the slightly weaker theorem stated just below.
\end{Rmk}
%
\begin{Thm}
Define
\[ \mathit{Opt}' \big( (\epsilon_{ij})_{(i,j)\in I\times J} \big)
\coloneqq \sup \Big\{ \big\{ \vec{X}_I : \vec{Y}_J \big\} \ ;\ \big(\forall (i,j) \in I\times J\big) \,
\big(\{ X_i:Y_j \}_* \leq \epsilon_{ij}\big) \Big\} ;\]
then for any closed cone $C$ of~$\RR^{N\times M}$
contained in~$(\RR_+^*)^{N\times M} \cup \{0\}$,
on~$C$, one has
\[ \mathit{Opt}'(\vec{\epsilon}) \stackrel{\vec\epsilon\longto\vec0}\sim
\VERT \boldepsilon \VERT .\]
\end{Thm}




\section{Machinery for using the tensorization theorems}%
\label{parMachineryForUsingTheTensorizationTheorems}

Up to now we stated the tensorization theorems in a rather `theoretical' form.
To apply these results to `concrete' situations, some additional techniques may be needed.
This section gives such techniques, which we will use later for the applications
of Chapter~\ref{parApplications}.

\begin{NOTA}
In this section, all the probability systems considered
will be endowed with their natural $\sigma$-metalgebras,
cf.\ Definition~\ref{def3128}.
To alleviate notation, I will give no names to these $\sigma$-metalgebras,
but will plainly denote ${\{X:Y\}}_*$ to mean
``the subjective decorrelation between~$X$ and~$Y$
seen from the natural $\sigma$-metalgebra of the underlying system''.
\end{NOTA}


\subsection{The `doubling-up' technique}\label{parTheDoublingUpTechnique}

\begin{Def}
For~$I$ a set and~$\mcal{R}$ a binary relation on~$I$,
$J_1,J_2\subset I$, we will say that ``$J_2$ is $\mcal{R}$-disjoint to~$J_1$''
if $(i,j)\in J_1\times J_2\ \Rightarrow\ i\notR j$.
\end{Def}

\begin{Lem}[`Doubling-up' lemma]\label{lem0306}
Let~$I$ be a (countable) set and let~$(X_i)_{i\in I}$ be random variables
such that for all~$i,j\in I$, $\{X_i:X_j\}_* \leq \epsilon_{ij}$
for a certain family of~$\epsilon_{ij} \in [0,1]$.

Let~$\mcal{R}$ be a binary relation on~$I$;
for~$i,j\in I$, denote $\epsilon^{\mcal{R}}_{ij} \coloneqq \1{i\notR j}\epsilon_{ij}$.

Define $\mbf{I} = I_1 \uplus I_2$ to be a disjoint union of two copies of~$I$;
denote by~$(i_1)_{i\in I}$, resp.~$(j_2)_{j\in I}$, the elements of~$I_1$, resp.~$I_2$.
Assume that the following holds for a certain $\epsilon\in[0,1]$:
``if $(Y_{i_{\kappa}})_{i_{\kappa}\in\mbf{I}}$ are random variables such that
$\forall i, j \in I\enskip\{Y_{i_1}:Y_{j_2}\}_{*}
\leq \epsilon^{\mcal{R}}_{ij}$, then $\{\vec{Y}_{I_1}:\vec{Y}_{I_2}\} \leq \epsilon$''.

Then for all~$J_1,J_2\subset I$ such that $J_2$ is $\mcal{R}$-disjoint to~$J_1$,
$\{\vec{X}_{J_1}:\vec{X}_{J_2}\}_{*} \leq \epsilon$.
\end{Lem}

\begin{Rmk}
The interest of Lemma~\ref{lem0306} is that, by proving
\emph{one} tensorization result on ${\{\vec{Y}_{I_1}:\vec{Y}_{I_2}\}}$,
one gets tensorization results
on \emph{all} the~$\{\vec{X}_{J_1}:\vec{X}_{J_2}\}$
for $J_2$ $\mcal{R}$-disjoint to~$J_1$.
\end{Rmk}

\begin{Xpl}\label{xpl1201}\strut
\begin{enumerate}
\item If you take for~$\mcal{R}$ the equality relation,
then Lemma~\ref{lem0306} gives a decorrelation result
for all disjoint~$J_1$ and~$J_2$.
\item\label{itm1202} If $I$ is equipped with a distance $\mathit{dist}$
and if you take $(i\mcal{R}j) \, \Leftrightarrow \, {\big(\mathit{dist}(i,j)<d_1\big)}$,
then you get a decorrelation result for all~$J_1$ and~$J_2$
such that $\mathit{dist}(J_1,J_2) {\geq d_1}$.
\end{enumerate}
\end{Xpl}

\begin{proof}
Assume that the hypotheses of the lemma hold
and let~$J_1,J_2\subset I$ with~$J_2$ $\mcal{R}$-disjoint to~$J_1$.
For~$i_{\kappa} \in \mbf{I}$, define
\[ Y_{i_{\kappa}} = \begin{cases} X_i &
\text{if\ \ $(\kappa=1\enskip\text{and}\enskip i\in J_1)$\ \ or\ \ $(\kappa=2\enskip\text{and}\enskip i\in J_2)$;} \\
\partial & \text{otherwise,} \end{cases} \]
for~$\partial$ some cemetery point in the range of none of the~$X_i$.
Since a constant variable is always independent of any variable,
the hypothesis ``$\{X_i:X_j\}_* \leq \epsilon_{ij}$'' for all~$i,j \in I$
implies that ${\{Y_{i_1}:Y_{j_2}\}_{*}} \leq \epsilon^{\mcal{R}}_{ij}$,
so, by the assumption of the lemma, ${\{\vec{Y}_{I_1}:\vec{Y}_{I_2}\}} \leq \epsilon$.
But $\vec{X}_{J_1}$ is $\vec{Y}_{I_1}$-measurable,
resp.\ $\vec{X}_{J_2}$ is $\vec{Y}_{I_2}$-measurable,
hence $\{\vec{X}_{J_1}:\vec{X}_{J_2}\} \leq \epsilon$.

Getting the subjective result w.r.t.~$*$
is just a variant of that reasoning, cf.\ \S~\ref{parSubjectiveResults}.
\end{proof}


\subsection{A practical result on~\texorpdfstring{$\ZZ^n$}{Z\textasciicircum n}}\label{parPractical}

In Chapter~\ref{parApplications}, the situations we will handle
shall always be of the following form:
%
\begin{Hyp}\label{hypZZn}
For some $n\in\NN^*$, the system is made of random variables $X_i$, $i\in\ZZ^n$,
which satisfy the condition
\[ \forall i,j\in\ZZ^n \qquad \{X_i:X_j\}_* \leq \epsilon(j-i) \]
for some symmetric function~$\epsilon \colon \ZZ^n \to [0,1]$.
\end{Hyp}

For systems satisfying Assumption~\ref{hypZZn},
one has the following practical synthetic result:
%
\begin{Lem}\label{pro8882}
Consider a norm $|\Bcdot|$ on~$\RR^n$,
the associated distance on the affine $\RR^n$ being denoted by~$\dist$.
Then for a system satisfying Assumption~\ref{hypZZn},
for all~$J_1,J_2\subset I$:
\[\label{for7041} \big\{ \vec{X}_{J_1} : \vec{X}_{J_2} \big\} \leq
\Big(\sum_{\mathclap{\substack{z\in\ZZ^n\\|z|\geq\dist(J_1,J_2)}}} \epsilon(z)\Big)
\wedge 1 .\]
\end{Lem}

\begin{proof}
To alleviate notation, denote $d \coloneqq \dist(J_1,J_2)$.
Applying Lemma~\ref{lem0306},
taking for ``$\mcal{R}$'' the relation ``be at distance $<d$''
(cf.\ Example~\ref{xpl1201}-\ref{itm1202}),
our goal becomes the following:
supposing $(Y_{i_\kappa})_{i_\kappa\in\ZZ^n_1\uplus\ZZ^n_2}$ are random variables
such that $\{Y_{i_1}:Y_{j_2}\}_* \leq \1{|j-i|\geq d}\epsilon(j-i)$,
we want to bound above $\{\vec{Y}_{\ZZ^n_1}:\vec{Y}_{\ZZ^n_2}\}$.

To do this we apply Theorem~\ref{thm6413},
and we get that $\{\vec{Y}_{\ZZ^n_1}:\vec{Y}_{\ZZ^n_2}\}$ is bounded by
$\VERT \boldepsilon \VERT \wedge 1$,
where $\boldepsilon$ is the following operator:
%
\[ \begin{array}{rrcl} \boldepsilon \colon & L^2(\ZZ) & \circlearrowleft & \\
& \big(g(j)\big)_{j\in\ZZ} & \mapsto &
\big( \sum_{j\in\ZZ} \1{|j-i|\geq d} \epsilon(j-i) g(j) \big)_{i\in\ZZ} .
\end{array} \]

To compute $\VERT \boldepsilon \VERT$, we split $\boldepsilon$
as $\sum_{z\in\ZZ} \1{|z|\geq d} \epsilon(z) M_z$, where
$M_z$ is the operator
\[ \begin{array}{rrcl} M_z \colon & L^2(\ZZ) & \circlearrowleft & \\
& \big(g(j)\big)_{j\in\ZZ} & \mapsto &
\big( g(i+z) \big)_{i\in\ZZ} .\end{array} \]
%
Obviously $\VERT M_z \VERT = 1$,
thus $\VERT \boldepsilon \VERT \leq \sum_{|z|\geq d} \epsilon(z)$
—actually there is even equality\hbox{—,}
which ends the proof of Lemma~\ref{pro8882}.
\end{proof}

\begin{Rmk}
Instead of Theorem~\ref{thm6413}, here we could have used Theorem~\ref{corZnZn},
which would yield a better result;
yet that would be very specific to~$\ZZ^n$ (cf.\ \S~\ref{parNonFlat}),
and the result would actually be almost equivalent to~(\ref{for7041})
(cf.\ \S~\ref{parAsymptoticOptimality}).
\end{Rmk}


\subsection{Avoiding the artificial phase transition}%
\label{parAvoidingThePhaseTransitionForTheToyModel}

Let us look again at Formula~(\ref{for7041}):
the ``$\wedge 1$'' in it is not really relevant
since a correlation level is \emph{always} bounded by~$1$.
In fact the situation is di\-cho\-to\-mic:
denoting $d \coloneqq \dist(I,J)$, either $\sum_{|z|\geq d} \epsilon(z)$ is $<1$ and then
(\ref{for7041}) is a true decorrelation result,
or it is $\geq1$ and then (\ref{for7041}) tells us actually nothing.
In other words, our result has a `phase transition'
depending on the relative values of~$\sum_{|z|\geq d} \epsilon(z)$ and~$1$,
similar to the phenomenon we discussed in Remark~\ref{rmkPhTr}.

However, as I pointed out in Footnote~\ref{ftn6129} on page~\pageref{ftn6129},
it is not clear whether the phase transition we are dealing with is a real phenomenon:
maybe it is rather an artifact due to Theorem~\ref{thm6413}'s bound's being non-optimal,
which could be avoided by a cleverer reasoning.
We are strengthened in that thought by observing that,
if $\sum_{z\in\ZZ^n} \epsilon(z) < \infty$,
then for~$d$ large enough one has
$\sum_{|z|\geq d} \epsilon(d) < 1$,
so that there is no phase transition for long distances;
why would a transition appear all of a sudden for short distances?

This subsection will show that, indeed, phase transitions can be avoided
in the situations we deal with.

\begin{Lem}\label{lem4067}
For a system satisfying Assumption~\ref{hypZZn}
with $\epsilon(z)<1$ as soon as $z\neq0$ and $\sum_{z\neq0} \epsilon (z) < \infty$,
there exists a constant $k<1$ such that, for all disjoint
$J_1,J_2\subset I$, one has ${\{\vec{X}_{J_1}:\vec{X}_{J_2}\}} \ab \leq k$.
\end{Lem}

\begin{proof}
As before, using Lemma~\ref{lem0306}
we have to bound above $\{\vec{Y}_{\ZZ^n_1}:\vec{Y}_{\ZZ^n_2}\}$
in the relevant doubled-up model.
Our plan to avoid the phase transition will consist in reducing to the `long distance' case.

For some $\ell \in \NN^*$,
we split $\ZZ^n_1$, resp.~$\ZZ^n_2$,
into a partition of~$\ell^n \eqqcolon N$ sublattices $Z_1^{(1)}, \ab \ldots, \ab Z_1^{(N)}$,
resp.~$Z_2^{(1)}, \ab \ldots, \ab Z_2^{(N)}$,
each lattice $Z_\kappa^{(u)}$ being of the form~$\ell \ZZ^n + z_u$ for some $z_u\in\ZZ^n \div \ell\ZZ^n$.
I claim two fundamental properties of these sublattices:

\begin{Clm}\label{clm8133}
For all~$u,v\in\{1,\ldots,N\}$,
\[\label{for8590} \big\{ \vec{Y}_{\ZZ_1^{(u)}} : \vec{Y}_{\ZZ_2^{(v)}} \big\}_*
\leq \sum_{z \equiv z_v-z_u} \!\! \1{z\neq0} \, \epsilon(z) .\]
\end{Clm}

\begin{proof} It is analogous to the proof of Lemma~\ref{pro8882}.
\end{proof}

\begin{Clm}\label{clm9816}
Provided $\ell$ is large enough,
the right-hand side of~(\ref{for8590}) is (strictly) less than $1$
for all the possible values of~$u,v$.
\end{Clm}

\begin{proof}
Denote $\zeta \coloneqq \sup_{z\neq0} \epsilon(z)$;
notice that our assumptions imply that $\zeta < 1$.
Since $\sum_{z\in\ZZ^n} \epsilon(z)$ converges,
there exists some $d_1<\infty$ such that $\sum_{|z|>d_1} \epsilon(z) < 1-\zeta$.
Now, denoting $d_0 \coloneqq \min\{ |z| \colon z\in\ZZ^n\setminus\{0\} \}$,
for $\ell > 2d_1 \div d_0$, for all~$u,v$ there is at most one $z$ congruent to
$z_v-z_u$ [mod.\ $\ell$] such that $|z| \leq d_1$,
whence the following uniform bound for the right-hand side of~(\ref{for8590}):
%
\[ \sum_{z\equiv z_v-z_u} \!\! \1{z\neq0} \, \epsilon(z)
\leq \sum_{|z| > d_1} \epsilon(z)
+ \underbrace{\sum_{\substack{|z|\leq d_1\\z\equiv z_v-z_u\\z\neq 0}} \epsilon(z)}%
_{\substack{\leq \zeta \text{ because}\\\text{the sum has at most}\\\text{one term, being $\leq \zeta$}}}
\leq \underbrace{\sum_{|z| > d_1} \epsilon(z)}_{< 1-\zeta} + \zeta < 1.\]
\end{proof}

Now, suppose $\ell$ large enough so that Cl\-aim~\ref{clm9816} works.
We apply \emph{simple} tensorization (Theorem~\ref{thm5750}) to the~$\vec{Y}_{\ZZ_2^{(v)}}$:
writing that $\vec{Y}_{\ZZ^n_2} = \big( \vec{Y}_{\ZZ_2^{(1)}}, \ldots, \vec{Y}_{\ZZ_2^{(N)}} \big)$,
we get that, for any~$u\in\{1,\ldots,N\}$,
\[\label{for8741} \big\{ \vec{Y}_{\ZZ_1^{(u)}} : \vec{Y}_{\ZZ^n_2} \big\}_*
\leq \sqrt{ 1 - \prod_{v=1}^{N}
\left( 1 - \big\{ \vec{Y}_{\ZZ_1^{(u)}}:\vec{Y}_{\ZZ_2^{(v)}} \big\}_*^2 \right) } < 1 .\]
%
Now we write $\vec{Y}_{\ZZ^n_1} = \big( \vec{Y}_{\ZZ_1^{(1)}}, \ldots, \vec{Y}_{\ZZ_1^{(N)}} \big)$
and we apply simple tensorization again\ab —this time to the~$\vec{Y}_{\ZZ_1^{(u)}}$ —to get:
\[\label{for1198}
\big\{ \vec{Y}_{\ZZ^n_1} : \vec{Y}_{\ZZ^n_2} \big\}
\leq \sqrt{ 1 - \prod_{u=1}^{N}
\left( 1 - \big\{ \vec{Y}_{\ZZ_1^{(u)}}:\vec{Y}_{\ZZ^n_2} \big\}_*^2 \right) } < 1 .\]
Bound~(\ref{for1198}) achieves our goal.
\end{proof}

\begin{Rmk}
With that proof, the way $k$ depends on~$\epsilon(\Bcdot)$ is rather complicated;
in particular, you cannot express $k$
as a function of only~$\sum_{z\neq0} \epsilon(z)$ and~$\sup_{z\neq0} \epsilon(z)$.
\end{Rmk}

\begin{Rmk}
In the case $n=1$, at first sight Lemma~\ref{lem4067} seems to contradict Theorem~\ref{t6657},
in which we told that Theorem~\ref{thm0119}, which does have a phase transition, was optimal.
The explanation for this paradox stands in the slight difference between the assumptions
of Lemma~\ref{lem4067} and Theorem~\ref{thm0119}:
while in Lemma~\ref{lem4067} we really imposed that $\{X_i:X_j\}_* \leq \epsilon\big(\mfrak{d}(i,j)\big)$,
with ``$*$'' denoting the \emph{full} natural $\sigma$-metalgebra of the system,
in Theorem~\ref{thm0119}—more precisely, in the version of Theorem~\ref{thm0119}
Theorem~\ref{t6657} proved to be optimal, which was the \emph{minimal} version of this theorem
(cf.\ \S~\ref{parMinimalHypotheses})—%
the conditions on subjective decorrelations were a bit looser.
That difference makes all the trick when one performs the steps of simple tensorization in the proof of Lemma~\ref{lem4067},
because these steps require subjective decorrelations w.r.t.\ the~$\vec{Y}_{I_\kappa^{(u)}}$,
which the sole assumptions of Theorem~\ref{t6657} do not provide.
\end{Rmk}


\subsection{Non-flat geometries}\label{parNonFlat}

It is natural to ask what we one can do
when the basic variables $X_i$ are not indexed by~$\ZZ^n$,
but by the vertices of a more general graph,
for instance a tree or a finitely generated group.
This shall occur indeed if the physical space one works in
exhibits some curvature%
—though Chapter~\ref{parApplications} will not handle such situations.

Actually for general graphs there are results analogous to those of
\S\S~\ref{parPractical} and~\ref{parAvoidingThePhaseTransitionForTheToyModel},
with similar (though more technical) proofs~\cite{perso}.
Here I will only give the statements of these results.

In this subsection the situation will be the following:
%
\begin{Hyp}\label{hypNF}
The system is made of random variables $(X_i)$ indexed by a (countable) set $I$.
There is a group $G$ acting transitively on~$I$,
and $I$ is endowed with a symmetric map $\mfrak{d} \colon {I\times I}\to\mfrak{D}$,
called the `abstract distance', which is preserved by the action of~$G$.
We assume that one has
\[ \forall i,j\in I \quad \{X_i:X_j\}_{*} \leq \epsilon\big(\mfrak{d}(i,j)\big) \]
for some function~$\epsilon \colon \mfrak{D} \to [0,1]$.
\end{Hyp}

\begin{Def}
For~$d\in\mfrak{D}$, we define $\valency(d) \coloneqq \# \{ i\in I \colon \mfrak{d}(o,i) = d \}$,
where the choice of~$o\in I$ does not matter.
\end{Def}

Then the analogous to Lemma~\ref{pro8882} is the
%
\begin{Lem}\label{pro8882+}
For~$\mfrak{D}'\subset\mfrak{D}$, for all~$J_1,J_2\subset I$ such that
$(i\in J_1,j\in J_2)\,\Rightarrow \mfrak{d}(i,j) \ab \in \nolinebreak[2] \mfrak{D}'$,
\[ \big\{\vec{X}_{J_1} : \vec{X}_{J_2}\big\} \leq
\Big( \sum_{d\in \mfrak{D}'} \valency(d)\,\epsilon(d) \Big) \wedge 1 .\]
\end{Lem}

The analogous of Lemma~\ref{lem4067} is the
%
\begin{Lem}\label{lem4067+}
Assume that Assumption~\ref{hypNF} is satisfied;
denoting by~$0$ the (common) value of the~$\mfrak{d}(i,i)$,
also assume that $\valency(0)=1$ and that $\epsilon(d) < 1$ as soon as $d \neq 0$.
Assume that $\sum_{d\in\mfrak{D}} \valency(d)\epsilon(d) < \infty$.

Moreover, assume that the action of~$G$ on~$I$ is \emph{profinite}
(cf.~\cite[Definition~1.1]{profinite}), i.e.\ that
there is a subset $\mcal{N} \subset \NN^*$
such that for each~$N\in\mcal{N}$, there is a subgroup $G_N \leq G$ such that:
%
\begin{ienumerate}
\item\label{itm4759} The action of~$G_N$ splits $I$ into exactly $N$ orbits
$I^{(1)},\ldots,I^{(N)}$;
\item\label{itm4518} $G_N$ is normal,
so that the partition of~$I$ into the~$I^{(u)}$ is stable by the action of~$G$;
\item\label{itm4750} Any two distinct points of~$I$ are ultimately separated
by the partitions induced by the~$G_N$, i.e.:
\[ \limsup_{\substack{N\in\mcal{N}\\N\longto\infty}} (G_{N}\cdot o) = \{o\} .\]
\end{ienumerate}

Then there exists a constant $k<1$ such that, for all disjoint
$J_1,J_2\subset I$, one has $\{\vec{X}_{J_1}:\vec{X}_{J_2}\} \leq k$.
\end{Lem}

\begin{Xpl}
For $I=\ZZ^n$ on which $G=\ZZ^n$ acts by translation,
equipped with the abstract distance $\mfrak{d}(x,y) = \{\pm(y-x)\}$,
the assumptions of Lemmas~\ref{pro8882+} and~\ref{lem4067+} are checked,
and these lemmas re-give resp.\ Lemmas~\ref{pro8882} and~\ref{lem4067}.
\end{Xpl}

\begin{Xpl}\label{x3339}
For $I$ the modular group $\mathit{PSL}_2(\ZZ)$ acting by left multiplication on itself,
equipped with its natural abstract distance (i.e., $\mfrak{d}(i,j) = \{i^{-1}j,j^{-1}i\}$),
the assumptions of Lemmas~\ref{pro8882+} and~\ref{lem4067+} are also checked%
—to see that the action of~$G$ on~$I$ is profinite,
take for the~$G_{N(\ell)}$ the principal congruence subgroups
$\Gamma(\ell)$ of the modular group~\cite{PSL2Z}.
%
Notice that $\mathit{PSL}_2(\ZZ)$ is an example of graph having negative curvature%
~\cite{Gromov}.
\end{Xpl}




\section{Appendix: Illustration of the proof of Theorem~\ref{thm6413}}%
\label{parIllustrationOfTheProofs}

\begin{NOTA}
This subsection is devised for the readers who would like to
understand better the proof of Theorem~\ref{thm6413}
by seeing how it works on a concrete example.
It only contains pedagogical material,
and thus can be skipped safely.
\end{NOTA}


\subsection{A Gaussian system of variables}\label{parX1X2Y}

In this illustration we take $N=2, M=1$ —since $M=1$, $Y_1$ will merely be denoted by~$Y$ —,
and we take $(X_1,X_2,Y)$ Gaussian (and centered),
whose law is described through a $3\times3$ matrix
\emph{via} writing that, for some standard Gaussian vector $(\xi_1,\xi_2,\xi_3) \in \RR^3$,
\[\label{for5012}
\begin{pmatrix} X_1\\X_2\\Y \end{pmatrix} =
\begin{pmatrix} \alpha_1 & \alpha_2 & \alpha_3 \\
\beta_1 & \beta_2 & \beta_3 \\ \omega_1 & \omega_2 & \omega_3 \end{pmatrix}
\begin{pmatrix} \xi_1\\\xi_2\\\xi_3 \end{pmatrix}.\]
We denote the matrix appearing in~(\ref{for5012}) by~$\mbf{M}$.
The rows of~$\mbf{M}$ will be denoted by~$\alpha,\beta,\gamma \in \RR^3$,
and $(\xi_1,\xi_2,\xi_3)$ will be denoted by~$\xi\in\RR^3$.
On~$\RR^3$ we will use the Euclidian scalar product ``$\cdot$'' and the associated norm ``$\|\Bcdot\|$''.

The advantage of this model is that,
by of the general properties of Gaussian vectors
(in particular Theorem~\ref{pro1857}),
all the quantities of interest are computable exactly.

First we compute the correlation levels:
by Theorem~\ref{pro1857},
\[ \{X_1:Y\} = \frac{|\alpha \cdot \omega|}{\|\alpha\| \, \|\omega\|} ,\]
similarly $\{X_2:Y\} = |\beta \cdot \omega| \div \|\beta\|\|\omega\|$; and
\[ \{\vec{X}:Y\} = \sqrt{1-\frac{|\omega\cdot(\alpha\vec{\times}\beta)|^2}{\|\omega\|^2\,\|\alpha\vec{\times}\beta\|^2}}, \]
where ``$\vec{\times}$'' denotes the cross product on $\RR^3$.
%
Concerning the conditional quantities,
denote by~$\beta^1$, resp.~$\omega^1$, the (orthogonal) projections of~$\beta$, resp.~$\omega$, on~$\RR\alpha$,
and~$\beta^*$, resp.~$\omega^*$, the projections of the same vectors on~$(\RR\alpha)^\perp$,
i.e.\ (assuming that $\alpha\neq0$):
\[ \beta^1 \coloneqq (\beta\cdot\alpha \div \|\alpha\|^2)\,\alpha , \quad
\omega^1 \coloneqq (\omega\cdot\alpha \div \|\alpha\|^2)\,\alpha ; \]
\[ \beta^* \coloneqq \beta - \beta^1 , \quad
\omega^* \coloneqq \omega - \omega^1 \]
(see Figure~\ref{fig-abo}).
%
\begin{figure}
\centering
%
%% Dessin pour \beta^1, \omega^1, \beta^*, \omega^*
%
\noindent\strut\hspace{1em}%
\begin{tikzpicture}[>=stealth,x=.5cm,y=.5cm,inner sep=1.5pt]
\coordinate (O) at (0,0);
\coordinate (A) at (0,5.5);
\coordinate (Bb) at (0,4.5);
\coordinate (Zb) at (0,2.5);
\coordinate (Bt) at (-2.2,-1.8);
\coordinate (Zt) at (5.6,-1.2);
\coordinate (B) at ($(Bb)+(Bt)$);
\coordinate (Z) at ($(Zb)+(Zt)$);
\draw (O) node[anchor=west]{$0$};

%% "Bb" signifie "beta_zéro", "Zt" signifie "omega_étoile", etc..

\draw[very thick,->] (O)--(A) node[anchor=east]{$\alpha$};
\draw[very thick,->] (O)--(B) node[anchor=east]{$\beta$};
\draw[very thick,->] (O)--(Z) node[anchor=west]{$\omega$};
\draw[semithick,->] (O)--(Bb) node[anchor=west]{$\beta^1$};
\draw[semithick,->] (O)--(Zb) node[anchor=east]{$\omega^1$};
\draw[semithick,->] (O)--(Bt) node[anchor=east]{$\beta^*$};
\draw[semithick,->] (O)--(Zt) node[anchor=west]{$\omega^*$};
\draw[very thin, dashed] (B)--(Bb);
\draw[very thin, dashed] (B)--(Bt);
\draw[very thin, dashed] (Z)--(Zb);
\draw[very thin, dashed] (Z)--(Zt);
\draw[thin] (-3.6,0.2)--(-6.3,-2.5) node[above right]{\scriptsize~~$(\RR\alpha)^\perp$} --(7.2,-2.5);

\rightangle{B}{Bb}{O}
\rightangle{Z}{Zb}{O}
\rightangle{B}{Bt}{O}
\rightangle{Z}{Zt}{O}
\end{tikzpicture}
%
\hfill
%
%% Dessin pour les autres vecteurs, relatifs au machin conditioné par Y.
%
\begin{tikzpicture}[>=stealth,x=.5cm,y=.5cm,inner sep=1.5pt]
\coordinate (O) at (0,0);
\coordinate (Z) at (0,5.4);
\coordinate (Bd) at (-2.2,1.1);
\coordinate (Bc) at (-2.7,-.9);
\coordinate (At) at ($2.2*(Bc)$);
\coordinate (Ab) at (0,3.2);
\coordinate (Bb) at (0,4.1);
\coordinate (Bt) at ($(Bc)+(Bd)$);
\coordinate (B) at ($(Bt)+(Bb)$);
\coordinate (A) at ($(At)+(Ab)$);
\draw (O) node[anchor=north west]{$0$};

%% "Bc" : "beta_chapeau" ; "Bd" : "beta_dague", etc..

\draw[very thick,->] (O)--(A) node[anchor=east]{$\alpha$};
\draw[very thick,->] (O)--(B) node[anchor=east]{$\beta$};
\draw[very thick,->] (O)--(Z) node[anchor=east]{$\omega$};
\draw[semithick,->] (O)--(Ab) node[anchor=west]{$\bar\alpha$};
\draw[semithick,->] (O)--(At) node[anchor=east]{$\tilde\alpha$};
\draw[semithick,->] (O)--(Bb) node[anchor=west]{$\bar\beta$};
\draw[semithick,->] (O)--(Bt) node[anchor=north east]{$\tilde\beta$};
\draw[semithick,->] (O)--(Bc) node[anchor=north west]{$\hat\beta$};
\draw[semithick,->] (O)--(Bd) node[anchor=east]{$\beta^\dag$};
\draw[very thin, dashed] (A)--(Ab);
\draw[very thin, dashed] (A)--(At);
\draw[very thin, dashed] (B)--(Bb);
\draw[very thin, dashed] (B)--(Bt);
\draw[very thin, dashed] (Bt)--(Bc);
\draw[very thin, dashed] (Bt)--(Bd);
\draw[thin] (-7.5,0.3)--(-10.5,-2.7) node[above right]{\scriptsize~~$(\RR\omega)^\perp$} --(1.5,-2.7);

\rightangle{B}{Bb}{O}
\rightangle{A}{Ab}{O}
\rightangle{B}{Bt}{O}
\rightangle{A}{At}{O}
\rightangle{Bt}{Bc}{O}
\rightangle{Bt}{Bd}{O}
\end{tikzpicture}%
\hspace{1em}\strut
%
\caption{Visual definitions of the vectors derived from $\alpha$, $\beta$ and~$\omega$:
the left drawing shows how to build $\beta^1,\ab \omega^1,\ab \beta^*,\ab \omega^*$;
the right drawing (with different values for~$\alpha,\ab \beta,\ab \omega$) explains
the construction of~$\protect\bar\alpha,\ab \protect\bar\beta,\ab \protect\tilde\alpha,\ab \protect\tilde\beta,\ab \protect\hat\beta,\ab \beta^\dag$.}%
\label{fig-abo}
%
\end{figure}%
%
Then one has
$\EE[X_2|X_1] = \beta^1_1\xi_1 + \beta^1_2\xi_2 + \beta^1_3\xi_3 = \beta^1\cdot\xi$,
resp.\ $\EE[Y|X_1] = \omega^1\cdot\xi$,
thus $X_2-\EE[X_2|X_1] = \beta^*\cdot\xi$,
resp. $Y-\EE[Y|X_1] = \omega^*\cdot\xi$.
%
As $(X_1,X_2,Y)$ is Gaussian,
the law of~$(X_2-\EE[X_2|X_1],\ab {Y-\EE[Y|X_1]})$
under~$\Pr[\Bcdot|\ab X_1=x]$ does not depend on the value of~$x$;
therefore we know all the conditional laws of~$(X_2,Y)$ under the~$\Pr[\Bcdot|X_1=x]$,
and for all these laws $\{X_2:Y\}$ is equal by Theorem~\ref{pro1857} to
$|\beta^*\cdot\omega^*| \div \|\beta^*\|\|\omega^*\|$, so in the end:
\[ \{ X_2:Y \}_{X_1} = \frac{|\beta^*\cdot\omega^*|}{\|\beta^*\|\,\|\omega^*\|} .\]
$\{ X_1:Y \}_{X_2}$ can be computed by a similar formula.

Now let us `dissect' the proof of Theorem~\ref{thm6413} on our example.
We take $f$ linear, namely
\[\label{forX1+X2} f(X_1,X_2) \coloneqq X_1 + X_2 ,\]
so that all the computations shall again be tractable exactly.

Let us start with computing the quantities linked to~$f^0$:
one has
\begin{eqnarray}
f^0 = f &=& (\alpha+\beta)\cdot\xi ;\\
f_1^0 = f^{\sigma(X_1)} &=& X_1 + (X_2)^{\sigma(X_1)} = (\alpha+\beta^1)\cdot\xi ;\\
f_2^0 = f - f^{\sigma(X_1)} &=& \beta^*\cdot\xi ,
\end{eqnarray}
whence respectively
\begin{eqnarray}
V = V^0 &=& \|\alpha+\beta\|^2 = \|\alpha\|^2 + \|\beta\|^2 + 2\alpha\cdot\beta ; \\
V_1^0 &=& \|\alpha+\beta^1\|^2 = \|\alpha\|^2 + 2\alpha\cdot\beta + \frac{(\alpha\cdot\beta)^2}{\|\alpha\|^2} ; \\
V_2^0 &=& \|\beta^*\|^2 = \|\beta\|^2 - \frac{(\alpha\cdot\beta)^2}{\|\alpha\|^2} .
\end{eqnarray}
By the way we check that, as claimed by Formula~(\ref{for3043}),
$V^0 = V_1^0 + V_2^0$.

Now we turn to the quantities linked to~$f^1$.
First we have to compute the conditional laws of~$(X_1,X_2)$
under the events ``$Y=y$''.
%
The technique is the same as for computing $\{X_2:Y\}_{X_1}$
a few lines above:
denoting by~$\bar{\alpha}$, resp.\ by~$\bar{\beta}$, the projections of~$\alpha$, resp.~$\beta$, on~$\RR\omega$,
and~$\tilde{\alpha}$, resp.~$\tilde{\beta}$, the projections of the same vectors on~$(\RR\omega^\perp)$, i.e.\ (see Figure~\ref{fig-abo})
\[ \bar{\alpha} \coloneqq (\alpha\cdot\omega \div \|\omega\|^2)\,\omega , \quad
\bar{\beta} \coloneqq (\beta\cdot\omega \div \|\omega\|^2) \,\omega ; \]
\[ \tilde{\alpha} \coloneqq \alpha - \bar{\alpha} , \quad
\tilde{\beta} \coloneqq \beta - \bar{\beta} ,\]
%
one has $\EE[X_1|Y] = \bar{\alpha}\cdot\xi$,
resp.\ $\EE[X_2|Y] = \bar{\beta}\cdot\xi$, thus
$X_1 - \EE[X_1|Y] = \tilde{\alpha}\cdot\xi$,
resp.\ $X_2 - \EE[X_2|Y] = \tilde{\beta}\cdot\xi$;
and $(X_1 - \EE[X_1|Y], X_2 - \EE[X_2|Y])$
has the same law under all the~$\Pr[\Bcdot|Y=y]$.
%
So we can compute the quantities linked to~$f^1$ in the same way
as we computed those linked to~$f^0$: denoting
\[ \hat{\beta} \coloneqq \frac{\tilde{\beta} \cdot \tilde{\alpha}}{\|\tilde{\alpha}\|^2} \, \tilde{\alpha} ; \]
\[\beta^\dag \coloneqq \tilde{\beta} - \hat{\beta} \]
(see Figure~\ref{fig-abo}), one finds
%
\begin{eqnarray}
f^1 &=& (\tilde{\alpha}+\tilde{\beta})\cdot\xi ;\\
f_1^1 &=& (\tilde{\alpha}+\hat{\beta})\cdot\xi ;\\
f_2^1 &=& \beta^\dag \cdot \xi ,
\end{eqnarray}
whence respectively:
\begin{eqnarray}
V^1 &=& \| \tilde{\alpha} + \tilde{\beta} \|^2 ;\\
V_1^1 &=& \|\tilde{\alpha}+\hat{\beta}\|^2 ;\\
V_2^1 &=& \|\beta^\dag\|^2.
\end{eqnarray}
%
As for~$f^0$, we check that $V^1 = V_1^1 + V_2^1$,
since $\tilde{\alpha} + \tilde{\beta}$ is the orthogonal sum of
$\tilde{\alpha} + \hat{\beta}$ and~$\beta^\dag$.
Moreover one always has $V^1 \leq V^0$, resp.\ $V_2^1 \leq V_2^0$:
the first inequality follows indeed from ${(\tilde{\alpha} + \tilde{\beta})}$'s being
the projection of~${(\alpha + \beta)}$ on~$(\RR\omega)^\perp$,
and the second one from $\beta^\dag$'s being
the projection of~$\beta^*$ on~${(\RR\omega+\RR\alpha)}^\perp$.
%
These inequalities are consistent with the following corollary of Claim~\ref{clm1908},
obtained by applying the claim conditionally to~$\mcal{G}_{j-1}$
with the role of ``$f$'' played by~$f^{j-1}$ and the role of ``$Y$'' played by~$Y_j$:
%
\begin{Pro}\label{clm3701}
For all~$1\leq j\leq M$, all~$0\leq i\leq N$,
\[ \sum_{i'>i} V_{i'}^j \leq \sum_{i'>i} V_{i'}^{j-1} .\]
\end{Pro}


\subsection{Numerical computations}\label{par411}

Now let us see a numerical example.
Our parameters will be chosen so that the function~$f$ defined by~(\ref{forX1+X2})
is optimal in the supremum~(\ref{for3878})
defining the Hilbertian correlation coefficient ${\{\vec{X}:Y\}}$;
other than that, the behaviour of our example will be generic:
\[ \mbf{M} = \begin{pmatrix} 4&1&1 \\ 1&4&1 \\ 1&1&4 \end{pmatrix} .\]
%
For that $\mbf{M}$ the calculations of the previous subsection give:

\begin{center}
\renewcommand\tabcolsep{5pt}
\noindent\begin{tabular}{|c||c|c|c|c|c|c|c|c|c|c|c|c|c|c|c|c|c|c|}
\hline
$\chi$  & $\alpha$ & $\beta$ & $\gamma$ & $\alpha\vec{\times}\beta$ &
$\beta^1$ & $\omega^1$ & $\beta^*$ & $\omega^*$ &
$\bar{\alpha}$ & $\bar{\beta}$ & $\tilde{\alpha}$ & $\tilde{\beta}$ & 
$\hat{\beta}$ & $\beta^\dag$ \\
\hline\hline
$\chi_1$ & $4$ & $1$ & $1$ & $-3$ & $2$  & $2$  & $-1$ & $-1$  & $1/2$ & $1/2$ & $7/2$ & $1/2$ & $7/6$ & $-2/3$ \\
\hline
$\chi_2$ & $1$ & $4$ & $1$ & $-3$ & $1/2$ & $1/2$ & $7/2$ & $1/2$ & $1/2$ & $1/2$ & $1/2$ & $7/2$ & $1/6$ & $10/3$ \\
\hline
$\chi_3$ & $1$ & $1$ & $4$ & $15$ & $1/2$ & $1/2$ & $1/2$ & $7/2$ & $2$  & $2$  & $-1$ & $-1$ & $-1/3$ & $-2/3$ \\
\hline
\end{tabular}
\end{center}

\noindent whence $\{X_1:Y\} = 1/2$ and $\{X_1:Y\}_{X_2} = 1/3$,
thus $\{X_1:Y\}_{\mcal{M}} = 1/2$;
and similarly $\{X_2:Y\}_{\mcal{M}} = 1/2$.

Then Theorem~\ref{thm5750} yields:
\[ \{ \vec{X} : Y \} \leq 1/\sqrt{2} = 0.707\ldots ,\]
and even, according to the refinements of \S~\ref{parMinimalHypotheses}:
\[ \{ \vec{X} : Y \} \leq \sqrt{13}/6 = 0.600\ldots ;\]
on the other hand, the true result is:
\[ \{ \vec{X} : Y \} = 1/\sqrt{3} = 0.577\ldots .\]
So here the bound~(\ref{for3009}) is (fortunately!) correct, and even rather sharp.


Now, as the proof of Theorem~\ref{thm6413} consists in studying the relations
between the~$V^j_i$, let us see what these quantities look like here.
One computes:
\[ \begin{pmatrix} V_1^0 & V_2^0 \\ V_1^1 & V_2^1 \end{pmatrix} =
\begin{pmatrix} 40\frac{1}{2} & 13\frac{1}{2} \\ 24 & 12 \end{pmatrix} .\]
As a first consequence, we can check the conclusions of Proposition~\ref{clm3701}:
$V_2^1 = 12 \leq V_2^0 = 13\frac{1}{2}$,
resp.\ $V_1^1+V_2^1 = 36 \leq V_1^0+V_2^0 = 54$.
%
Going further, we check the conclusions of Claim~\ref{clm1959},
which forbids the differences~$V_2^0 - V_2^1$ and~$(V_1^0+V_2^0) - (V_1^1+V_2^1)$
to be too large: for the first difference, one has
$V_2^0 - V_2^1 = 1\frac{1}{2} \leq \epsilon_2^2 V_2^0 = 3\frac{3}{8}$%
\footnote{According to \S~\ref{parMinimalHypotheses},
one can replace $\epsilon_2 = 1/2$ by~$\{X_2:Y\}_{X_1} = 1/3$
in this inequality. Then the inequality even becomes an equality:
this is linked to the optimality of certain tensorization results for Gaussian variables,
cf.\ \S~\ref{parOptimality}.}, and for the second one,
$(V_1^0+V_2^0) - (V_1^1+V_2^1) = 18 \leq
\big( \epsilon_1 V_1^0 + \sqrt{V_2^0-V_2^1} \big)^2 = 19.419\dots$.


\subsection{Some traps to avoid}

To finish with this appendix,
I would like to comment on what is true or not about the~$V^j_i$ in general situations.
Proposition~\ref{clm3701} pointed out that for all~$\hat\imath \in \{0,\ldots,N\}$,
$\sum_{i>\hat\imath} V^j_{i}$ is a nonincreasing function of~$j$;
in particular, when one looks at the table of the~$V^j_i$,
the last term ($\hat\imath=N-1$), resp.\ the total ($\hat\imath=0$) of line $j$
can only decrease.
%
Moreover, if in some line $j$ all the~$V^j_i$ are zero from some position $\hat\imath+1$,
then this property remains true in all the lower lines $j'>j$.
That can be explained very simply, since
saying that all the~$V^j_i$ are zero from position $\hat\imath+1$
means indeed that $f$ is $(\mcal{G}_j\vee\mcal{F}_{\hat\imath})$-measurable,
hence \emph{a fortiori} $(\mcal{G}_{j'}\vee\mcal{F}_{\hat\imath})$-measurable.
%
The following example, in which $f$ turns out to be $2X_1$, illustrates this phenomenon:
\[ \mbf{M} = \begin{pmatrix} 1&0&0 \\ 1&0&0 \\ 1&0&1 \end{pmatrix}
\qquad \Rightarrow \qquad
\begin{pmatrix} V_1^0 & V_2^0 \\ V_1^1 & V_2^1 \end{pmatrix} =
\begin{pmatrix} 4 & 0 \\ 2 & 0 \end{pmatrix} .\]

However, keep careful: almost anything else
you would like to say about the table of the~$V_i^j$ would be false!
In particular, for $i<N$, $V_i^j$ is not a nonincreasing function of~$j$ in general;
it is not even true that $V_i^j = 0 \ \Rightarrow V_i^{j'>j} = 0$,
as shown by the following example:
\[\label{for9186a} \mbf{M} = \left(\begin{array}{ccc} 1&0&0 \\ \!\!-1&1&0 \\ 1&1&1 \end{array}\right)
\qquad \Rightarrow \qquad
\begin{pmatrix} V_1^0 & V_2^0 \\ V_1^1 & V_2^1 \end{pmatrix} =
\begin{pmatrix} 0 & 1 \\ 1/6 & 1/2 \end{pmatrix} .\]

It is not true either that,
if $V^j$ remains unchanged from one line to another
(that is, the total of the~$V^j_i$ remains unchanged),
then all the~$V^j_i$ are unchanged.
In fact, that $V^{j+1}$ is equal to~$V^j$ means that, conditionally to~$\mcal{G}_j$,
$f^j$ is centered w.r.t.~$Y_{j+1}$, and then $f^{j+1} = f^j$.
However, the way $f^{j+1}$ decomposes into a sum of~$f^{j+1}_i$
may be different to the way $f^j$ decomposed into a sum of~$f^j_i$,
because conditioning w.r.t.~$Y_{j+1}$ may make the law of the~$X_i$ change!
That is what happens in the following example:
%
\[\label{for9186b} \mbf{M} = \begin{pmatrix} 1&0&1 \\ 0&1&\!\!-1 \\ 0&0&1 \end{pmatrix}
\qquad \Rightarrow \qquad
\begin{pmatrix} V_1^0 & V_2^0 \\ V_1^1 & V_2^1 \end{pmatrix} =
\begin{pmatrix} 1/2 & 3/2 \\ 1 & 1 \end{pmatrix} .\]




\section{Appendix: A corollary of the Perron--Fro\-be\-nius theorem}%
\label{parACorollaryOfThePerronFrobeniusTheorem}

In this appendix I handle a lemma used
in the proof of Theorem~\ref{thm6413}.
We are working on the vector space $\RR^N$ for some $N>0$;
a vector or a matrix is said to be ${>0}$ if all its entries are positive,
resp.~${\geq 0}$ if all its entries are nonnegative.
Then the \emph{Perron--Frobenius theorem}~\cite[Theorem~8.3.1]{HornJohnson} states that
if a square matrix $A$ is ${\geq 0}$, then $A$ has some ${\geq 0}$ eigenvector
for the eigenvalue $\rho(A)$. Our goal here is prove the following corollary:
%
\begin{Lem}\label{l3542}
Let~$A \geq 0$ be a square matrix,
then:
\[\label{for6491} \inf \big\{ \lambda\geq 0 \colon (\exists u > 0)(Au \leq \lambda u) \big\} = \rho(A) .\]
\end{Lem}

\begin{proof}
We prove separately each way of the equality.
Let us begin with way~``$\leq$''.
Let~$v \geq 0$ be some eigenvector of~$A$ for the eigenvalue $\rho(A)$.
If $v > 0$, then the value $\lambda = \rho(A)$ checks the condition in the infimum and we are done.
Otherwise if $v \ngtr 0$, up to a permutation of indices it has the form
$(0,\ldots,0,v'_{n+1},\ldots,v'_N)$ with $0<n<N$ and all the~$v'_i$ positive.
Reasoning by induction,
assume that we have proved the way~``$\leq$'' of the lemma
for all~$n<N$.
Then the form of the eigenvector~$v$ forces~$A$ to write blockwise
\[ A = \left( \begin{array}{cc} \tilde{A} & 0 \\ * & * \end{array} \right) \]
with $\RR^{n\times n} \ni \tilde{A} \geq 0$.
I claim that $\rho(\tilde{A}) \leq \rho(A)$, since
if $\tilde{v}$ is an eigenvector of~$\tilde{A}$ for the eigenvalue~$\rho(\tilde{A})$,
then for~$t \geq 0$
\[ A^{t} \big(\tilde{v},0,\ldots,0\big) = \big(\rho(\tilde{A})^{t}\,\tilde{v},*,\ldots,*\big) ,\]
so
\[\label{for5711}
\limsup_{t\longto\infty} \rho(\tilde{A})^{-t} \big|A^{t}(\tilde{v},0,\ldots,0)\big| > 0 \footnotemark \]
\footnotetext{Equation~(\ref{for5711}) is meaningless if $\rho(\tilde{A})=0$,
but in this case there is nothing to prove.}%
and consequently $\rho(A) \geq \rho(\tilde{A})$.
Now let~$\epsilon > 0$.
By induction hypothesis
there exists some $\RR^n \ni w \ab > 0$ such
that $\tilde{A}w \leq {(\rho(\tilde{A})+\epsilon) w}$.
Thus for~$\eta>0$, $\RR^N \ni (\eta w,v') > 0$ and
\[
A(\eta w,v') = \big( \eta \tilde{A} w,\rho(A)v'+O(\eta) \big)
\leq \big( \eta (\rho(\tilde{A})+\epsilon) w, \rho(A)v' + O(\eta) \big)
\stackrel{\eta \searrow 0}{\leq} \big(\rho(A)+\epsilon\big) (\eta w, v') .
\]
So $(\rho(A)+\epsilon)$ checks the condition in the right-hand side of the infimum,
which ends the proof of the way~``$\leq$'' of~(\ref{for6491}).

For the way~``$\geq$'', consider any~$\RR^N \ni u > 0$
and let again $v \geq 0$ be some eigenvector of~$A$
for the eigenvalue $\rho(A)$.
Then there exists a (unique) $\beta \geq 0$
such that $u-\beta v \geq 0$ but $u - \beta v \ngtr 0$.
For this $\beta$, one of the entries of~$\beta v$ and~$u$
is the same, say $\beta v_{i_0} = u_{i_0}$.
So if $\lambda < \rho(A)$,
\[ \lambda u_{i_0} < \rho(A) u_{i_0} = \rho(A) \beta v_{i_0} = \big(A(\beta v)\big)_{i_0}
\leq \big(A(\beta v)\big)_{i_0} + \big(A(u-\beta v)\big)_{i_0}
= (Au)_{i_0} ,\]
thus $Au \nleqslant \lambda u$.
That relation being true for any~$u>0$, $\lambda$ does not check the condition in the infimum,
which proves the way~``$\geq$'' of~(\ref{for6491}).
\end{proof}




\section{Appendix: A geometric consequence of results on correlations}\label{par3lines}

As I pointed out in Remark~\ref{rmk>3lines}, for Gaussian vectors
Hilbertian correlations can be interpreted in terms of Euclidian spaces.
In this appendix I will present a funny corollary of Lemma~\ref{l6791}
following from this interpretation.
Actually that result itself is more or less a pretext:
the real goal of this appendix is in fact to show in an eloquent way
the geometric meaning of maximal correlations
and the Hilbertian frame that underlies them.

First we need some vocabulary about Euclidian spaces:
%
\begin{Def}\strut
\begin{enumerate}
\item For~$L_1,L_2$ two vector lines in the Euclidian space $\RR^2$,
or more generally in any Hilbert space,
we call \emph{geometric angle} between~$L_1$ and~$L_2$,
denoted by~$\widehat{L_1L_2}$, their ``angle'' in the elementary sense:
for arbitrary $\vec{a}\in L_1\setminus\{0\}, \vec{b}\in L_2\setminus\{0\}$,
\[\widehat{L_1L_2}= \arccos \frac{|\vec{a}\cdot\vec{b}|}{\|\vec{a}\|\,\|\vec{b}\|} \in [0,\pi/2].\]
%
\item For~$L_1,L_2$ and~$L_3\neq L_1,L_2$ three vector lines in the Euclidian space $\RR^3$
(or any Hilbert space),
we call \emph{apparent angle between~$L_1$ and~$L_2$ seen from $L_3$}
the geometric angle that an observer located somewhere on~$L_3\setminus\{0\}$
would have the impression, due to perspective, that $L_1$ and~$L_2$ make
(see Figure~\ref{fig-3lines}): technically,
it is the geometric angle $\widehat{L'_1L'_2}$,
where $L_1'$ and~$L_2'$ are the respective orthogonal projections of~$L_1$ and~$L_2$
onto the plane $(L_3)^\perp$.
\end{enumerate}
\end{Def}

Then one has the following corollary of Lemma~\ref{l6791}:
%
\begin{Thm}\label{c8818}
Let~$L_1,L_2,L_3$ be three distinct vector lines of~$\RR^3$.
Denote $\widehat{A} \coloneqq \widehat{L_2L_3}, \widehat{B} \coloneqq \widehat{L_3L_1},
\widehat{\Omega} \coloneqq \widehat{L_1L_2}$,
and denote by~$\widehat{A}'$ the apparent angle between~$L_2$ and~$L_3$ seen from $L_1$,
resp.~$\widehat{B}'$ the apparent angle between~$L_3$ and~$L_1$ seen from $L_2$,
etc..
Then the relative order of~$\widehat{A}$ and~$\widehat{A}'$ is the same
as the relative order of~$\widehat{B}$ and~$\widehat{B}'$
and as the relative order of~$\widehat{\Omega}$ and~$\widehat{\Omega}'$, i.e.,
``$\widehat{A}'<\widehat{A}$'' (resp.\ ``$\widehat{A}'=\widehat{A}$'',
resp.\ ``$\widehat{A}'>\widehat{A}$'')
is equivalent to ``$\widehat{B}'<\widehat{B}$'' (resp.\ ``$\widehat{B}'=\widehat{B}$'',
resp.\ ``$\widehat{B}'>\widehat{B}$''),
etc..
\end{Thm}

\begin{Rmk}
I found Theorem~\ref{c8818} by chance, one day that I was looking for a situation
where one would have $\widehat{B}' > \widehat{B}$ but $\widehat{A}' < \widehat{A}$,
in order to build a `nice' example for \S~\ref{par411}.
I thought that such a situation would be generic,
but after having looked for it without success,
I realized that it was actually impossible, and that the explanation had a simple
interpretation in terms of correlations.
\end{Rmk}

\begin{figure}
\centering
\includegraphics[width=\textwidth]{quatredes}
\caption{This figure shows four different views of the same $3$-dimensional object.
What interests us actually is only the three concurrent lines $L_1,L_2,L_3$,
but we added a die
centered at their point of concurrency to see depth better on the pictures
{\footnotesize[Recall that on a die, the total number of points on two opposite faces is always $7$.]}.
%
On the top picture, the die is shown in generic position. We represent the angles~%
$\widehat{A}$, $\widehat{B}$ and~$\widehat\Omega$; these angles are $3$-dimensional angles,
which we underline by drawing them with double strokes.
%
On each of the bottom pictures, the die is viewed from the direction of one of the lines
(from the left to the right, $L_3$, $L_1$ and~$L_2$),
so that this line appears completely foreshortened.
We represent the angles $\widehat\Omega'$, $\widehat{A}'$ and~$\widehat{B}'$
made by the two other lines {\it as they appear on the drawing};
we underline that these angles are $2$-dimensional by drawing them with simple strokes
{\footnotesize[Note that in the case of~$\widehat{B}'$,
the angular sector representing $\widehat{B}'$
is not the projection of the angular sector representing $\widehat{B}$, but its supplementary%
—otherwise $\widehat{B}'$ would be greater than $\pi/2$, which would contradict
our `geometric' definition of angles.]}.
On this example one has
$\widehat{A}\simeq58\,^{\circ}, \widehat{B}\simeq71\,^{\circ}, \widehat\Omega\simeq15\,^{\circ}$ and $\widehat{A}'\simeq30\,^{\circ}, \ab \widehat{B}'\simeq34\,^{\circ}, \ab 
\widehat\Omega'\simeq9\,^{\circ}$; so, perspective makes angles appear
smaller than they are really for all three pairs of lines,
which is in accordance with Theorem~\ref{c8818}.}
\label{fig-3lines}
\end{figure}

\begin{proof}
Fix three arbitrary nonzero vectors $\alpha, \ab \beta, \ab \omega$ of resp.~$L_1, \ab L_2, \ab L_3$;
and consider the Gaussian system~(\ref{for5012}) of \S~\ref{parIllustrationOfTheProofs}
for these vectors.
Then the correlation coefficients between~$X_1$, $X_2$ and~$Y$
can be interpreted as angles between~$L_1$, $L_2$ and~$L_3$;
more precisely, one has the following correspondance:
%
\begin{Pro}\label{procscs}\strut
\begin{ienumerate}
\item\label{itmCos2}
$\{X_1:X_2\}$ is the cosine of the geometric angle between~$L_1$ and~$L_2$;
\item\label{itmCos3}
$\{X_1:X_2\}_Y$ is the cosine of the apparent angle between~$L_1$ and~$L_2$ seen from $L_3$.
\end{ienumerate}
\end{Pro}
%
\begin{proof}
(\ref{itmCos2}) is nothing but the Euclidian interpretation of Theorem~\ref{pro1857}.
(\ref{itmCos3}) follows from the fact that,
in the vector space spanned by jointly Gaussian real random variables,
conditional expectation corresponds to orthogonal projection
and independence corresponds to orthogonality.
\end{proof}

By Proprosition~\ref{procscs}, in our situation Lemma~\ref{l6791} gives:
\[\label{for0803a} \big\{(X_1,X_2):Y\big\} = \sqrt{1 - \sin^2\widehat{B} \, \sin^2\widehat{A}'} .\]
%
Obviously the roles of~$X_1$ and~$X_2$ can be interchanged in the above argument,
yielding:
\[\label{for0803b} \big\{(X_2,X_1):Y\big\} = \sqrt{1 - \sin^2\widehat{A} \, \sin^2\widehat{B}'} .\]
But $(X_1,X_2)$ and~$(X_2,X_1)$ generate the same $\sigma$-algebra,
so ${\{(X_1,X_2):Y\}} \ab = {\{(X_2,X_1):Y\}}$,
and thus, comparing~(\ref{for0803a}) and~(\ref{for0803b}):
\[ \frac{\sin\widehat{A}'}{\sin\widehat{A}} = \frac{\sin\widehat{B'}}{\sin\widehat{B}} .\]
This implies in particular that
$\sin\widehat{A},\sin\widehat{A}'$ and~$\sin\widehat{B},\sin\widehat{B}'$
have the same relative order, so also do $\widehat{A},\widehat{A}'$ and~$\widehat{B},\widehat{B}'$.
A cyclic permutation of~$L_1$, $L_2$ and~$L_3$
shows that the result is still valid for~$\widehat{\Omega},\ab \widehat{\Omega}'$.\linebreak[1]\strut
\end{proof}

\begin{Xpl}
See Figure~\ref{fig-3lines} on page~\pageref{fig-3lines}.
\end{Xpl}

